% !TeX document-id = {f4f6165e-7e15-441d-9a85-fe79757ad5fa}
%======================================================================
%  BACHELOR'S THESIS
%  Ruhr-Universität Bochum — IT-Sicherheit; carried out at MPI-SP
%  Layout: RUB thesis template (bossthesis.cls, bosscoverpage.sty, boss/),
%  adapted for MPI-SP; see the head of bossthesis.cls.
%
%  COMPILE:  latexmk (see latexmkrc): pdflatex -> biber -> pdflatex ...
%            Only the passes an edit actually needs are run.
%            (TeXstudio: F5 / F6 use latexmk via the comments below)
%
% !TeX TXS-program:compile = txs:///latexmk
% !TeX TXS-program:bibliography = txs:///biber
%
%  EDIT config/metadata.tex first. Then write in chapters/.
%======================================================================
\documentclass[twoside,its]{bossthesis}
% options of the class (see bossthesis.cls): [privacy] hides date and place
% of birth; [earlydraft] for quick printouts; [germanthesis] for German.


%----------------------------------------------------------------------
%  Metadata — EDIT THIS FILE, nothing else in config/
%  thesis.tex passes these to the bossthesis class (cover page, title
%  page, declaration). scripts/build_review_docx.py reads the
%  \newcommand lines for the review copy. Open points are inline todos:
%  they sit in the title page and the declaration box, where margin notes
%  cannot go.
%----------------------------------------------------------------------
\newcommand{\thesistitle}{Understanding Privacy Harm Communication in Health Tracking Apps}
\newcommand{\thesiskind}{Bachelor's Thesis in IT Security}
\newcommand{\thesiscitetype}{Bachelor's Thesis}
\newcommand{\authorname}{Hamza Sghaier}
\newcommand{\birthdate}{9. of June 2000}
\newcommand{\birthplacename}{Mahdia, Tunisia}
\newcommand{\advisorname}{Dr.~Kavous Salehzadeh Niksirat}
\newcommand{\examinerone}{Prof.~Dr.~Carmela Troncoso}
\newcommand{\examinertwo}{\todo[inline, inlinewidth=5cm]{Second examiner.}}
\newcommand{\thechair}{Max Planck Institute for Security and Privacy (MPI-SP)}
\newcommand{\thechairaddress}{Max Planck Institute for Security and Privacy\\Universitätsstraße 140 $\cdot$ 44799 Bochum $\cdot$ Germany}
\newcommand{\startdate}{\todo[inline]{Start date set by the examination office.}}
\newcommand{\submissiondate}{\todo[inline]{Submission date.}}

% AI tools named in the declaration box. Drafted from what the repository
% records; the author must check, correct and complete every line.
\newcommand{\aitoollist}{%
\begin{itemize}
 % \item \textit{Claude (Anthropic), used through Claude Code}: writing and
 %   debugging the collection, ranking and checking scripts in
 %   \texttt{scripts/} and the generated tables; preparing the eligibility
 %   screening decisions of the app sample, which the author checks
 %   (Section~\ref{sec:results-sample}); revising and drafting passages of
 %   Chapters~4 and~5 and Appendices~C and~E.
 %   \todo[inline]{Check and complete; name the sections.}
 % \item \textit{OpenAI Codex}: \todo[inline]{The repository contains a Codex
 %   configuration; state whether and for what it was used.}
 % \item \textit{Methods digest of the literature}: \todo[inline]{Name the tool
 %   that produced the cards and summaries in 01\_Literature/AI Literature
 %   Analysis, and its use.}
 % \item \textit{Language model of the coding pipeline}: part of the method,
 %   described in Chapter~\ref{ch:methodology} and Appendix~\ref{app:prompts}.
\item \todo[inline]{AI Usage summary.}
\end{itemize}}



%----------------------------------------------------------------------
%  Packages added to the bossthesis class (RUB template)
%  The class sets fonts, page layout, headings, page styles, captions,
%  hyperref, cleveref, babel and biblatex; this file adds what the thesis
%  needs beyond that and does not load those packages again.
%----------------------------------------------------------------------
\setlength{\emergencystretch}{1.5em}
% Pages may end short instead of stretching the space between paragraphs and
% list items (two-sided report pages are otherwise filled to full height).
\raggedbottom
\hyphenation{Kna-ben-hans}

% --- Colours ----------------------------------------------------------
% Restrained accents for diagrams, source listings and linked references.
% Body text, headings and table rules remain black.
\usepackage{xcolor}
\definecolor{primary}{HTML}{233F5B}
\definecolor{accent}{HTML}{37665D}
\colorlet{accentdark}{accent}
\definecolor{muted}{HTML}{535B63}
\definecolor{linkblue}{HTML}{496D91}

% --- Figures, tables, code ------------------------------------------
% graphicx, booktabs, float, tikz, listings and csquotes come with the class.
\graphicspath{{figures/}}
\usepackage{array}
\usepackage{colortbl}
\arrayrulecolor{black}
\setlength{\heavyrulewidth}{0.8pt}
\setlength{\lightrulewidth}{0.4pt}
\renewcommand{\arraystretch}{1.12}
\usetikzlibrary{arrows.meta}
\newcolumntype{L}[1]{>{\raggedright\arraybackslash}p{#1}}
\newcommand{\rowrule}{\cmidrule[0.3pt]{2-4}}
% Figure descriptions (from the proposal's ACM \Description) are not printed;
% they stay in the source for the review copy, which shows TikZ figures as text.
\newcommand{\Description}[1]{}
\usepackage{tabularx}
\usepackage{longtable}
\usepackage{pdflscape} % wide appendix tables on landscape pages
\colorlet{devshade}{primary!12} % rows of a developer with several apps
\captionsetup[table]{position=top, skip=6pt}
\captionsetup[figure]{skip=8pt}
\usepackage{enumitem}
\setlist{itemsep=0.2\baselineskip, topsep=0.4\baselineskip}
\setlist[itemize]{label=\small\textbullet}

% --- Links -----------------------------------------------------------
% The class loads hyperref and cleveref; links keep the thesis's blue.
\hypersetup{
  colorlinks = true,
  allcolors  = linkblue
}

% --- Bibliography ----------------------------------------------------
% The class loads biblatex (Biber, numeric-comp); restrained blue for
% linked citation numbers.
\addbibresource{../_shared/references.bib}
\renewcommand*{\bibfont}{\small}
\setlength{\bibitemsep}{0.35\baselineskip}
% Long URLs: break at letters and digits and set the list ragged right, so
% that lines are not stretched.
\setcounter{biburlnumpenalty}{9000}
\setcounter{biburllcpenalty}{7000}
\setcounter{biburlucpenalty}{8000}
\appto{\bibsetup}{\raggedright}
\DeclareFieldFormat{labelnumberwidth}{\textcolor{linkblue}{\mkbibbrackets{#1}}}

% --- Listings style --------------------------------------------------
% Subtle background and syntax accents; bold keywords also work in greyscale.
\definecolor{codebg}{HTML}{F5F6F7}
\colorlet{coderule}{primary!45}
\colorlet{codekeyword}{primary}
\definecolor{codebuiltin}{HTML}{665078}
\colorlet{codestring}{accent}
\colorlet{codecomment}{muted}
\lstdefinelanguage{json}{
  morestring = [b]",
  morekeywords = {true,false,null},
  sensitive = true
}
\lstdefinestyle{thesiscode}{
  basicstyle      = \linespread{1}\ttfamily\fontsize{8}{10}\selectfont,
  identifierstyle = {}, % the class sets bold blue identifiers
  keywordstyle    = \color{codekeyword}\bfseries,
  keywordstyle    = [2]\color{codebuiltin},
  stringstyle     = \color{codestring},
  commentstyle    = \color{codecomment},
  numberstyle     = \tiny\color{codecomment},
  numbers         = left,
  % Numbers stand outside the left rule with 8pt of clear space before it
  % (absolute units: listings sets em-lengths in the tiny number font).
  numbersep       = 12pt,
  stepnumber      = 1,
  frame           = leftline,
  rulecolor       = \color{coderule},
  backgroundcolor = \color{codebg},
  xleftmargin     = 26pt,
  framexleftmargin = 4pt,
  breaklines      = true,
  breakindent     = 1.2em,
  prebreak        = {},
  postbreak       = \mbox{\textcolor{codecomment}{$\hookrightarrow$}\space},
  columns         = fullflexible,
  keepspaces      = true,
  showstringspaces = false,
  tabsize         = 4,
  aboveskip       = 1ex,
  belowskip       = 1.5ex,
  captionpos      = t
}
\lstdefinestyle{pythoncode}{
  style    = thesiscode,
  language = Python,
  morekeywords = [2]{open,print,len,enumerate,any,dict,str,isinstance,set,list}
}
\lstdefinestyle{jsoncode}{style = thesiscode, language = json}
\lstset{style=thesiscode}

% --- Helpers ---------------------------------------------------------
% Revision marks are the template's todonotes: \todo{...} in the margin,
% \todo[inline]{...} where a margin note cannot go (tables, boxes, the
% title page). Quiet style: pale amber fill, thin amber border and leader
% line, small text, no shadow; the amber matches the orange \highlight
% marks. (size takes a size command only: todonotes also uses it inside
% \csname, so \sffamily there is an error.)
\definecolor{todofill}{HTML}{FFF5E1}
\definecolor{todoline}{HTML}{D9952E}
\definecolor{todotext}{HTML}{4A3410}
\setuptodonotes{color=todofill, bordercolor=todoline, linecolor=todoline,
  textcolor=todotext, size=\scriptsize}
\tikzset{notestyleraw/.append style={line width=0.4pt, rounded corners=1.5pt,
  inner sep=0.7ex}}
\usepackage{xpatch}
\makeatletter
% todonotes redraws the leader line "thick" for every note; make it thin.
\xpatchcmd{\@todo}{thick,}{line width=0.4pt,}{}{%
  \PackageWarning{preamble}{todonotes leader line not restyled}}
\makeatother
\usepackage{soul}

\newif\ifshowhighlights
\showhighlightstrue % Change to \showhighlightsfalse to hide highlights

% A todonotes leader line is drawn from the note's anchor along one text line
% only. A highlighted passage that is followed directly by \todo therefore
% gets a thin bottom edge in the note's colour on every line it spans, which
% runs into the leader line; a highlight without a note stays plain.
% \texthlul draws the highlight and the edge in one soul pass (soul cannot
% nest \hl and \ul); it uses soul's internals (\SOUL@...) as of TeX Live 2026.
\makeatletter
\newdimen\hl@lnht
\newdimen\hl@lnthick \hl@lnthick=0.7pt
\def\hl@ulpreamble{\SOUL@hlpreamble
  \hl@lnht=-\SOUL@uldp \advance\hl@lnht by \hl@lnthick}
\def\hl@ulleaders{\xleaders\hbox to 1pt{%
  \rlap{\vrule width 1.2pt height \SOUL@ulht depth \SOUL@uldp}%
  \rlap{\textcolor{todoline}{\vrule width 1.2pt height \hl@lnht depth \SOUL@uldp}}\hss}}
\DeclareRobustCommand*\texthlul{\SOUL@hlsetup
  \let\SOUL@preamble\hl@ulpreamble \let\SOUL@ulleaders\hl@ulleaders \SOUL@}
\newcommand{\hl@plain}[2]{{\sethlcolor{#1}\hl{#2}}}
\newcommand{\hl@marked}[2]{{\sethlcolor{#1}\texthlul{#2}}}
\newcommand{\hl@off}[2]{#2}
\newcommand{\hl@on}[2]{\@ifnextchar\todo{\hl@marked{#1}{#2}}{\hl@plain{#1}{#2}}}
\newcommand{\highlight}[2][yellow]{%
  \ifshowhighlights
    \expandafter\hl@on % \expandafter closes the conditional before the look-ahead
  \else
    \expandafter\hl@off
  \fi
  {#1}{#2}}
\makeatother



\author{\authorname}
\title{\thesistitle}
\thesistype{\thesiskind}
\thesiscite{\thesiscitetype}
\birthday{\birthdate}
\birthplace{\birthplacename}
\thesisstart{\startdate}
\thesisend{\submissiondate}
\institute{\thechair}
\instituteaddress{\thechairaddress}
\advisors{& \textbf{\advisorname}\\}
\examiners{First examiner: & \textbf{\examinerone}\\
           Second examiner: & \textbf{\examinertwo}\\}
\aitools{\aitoollist}

%======================================================================
\begin{document}

%---------------------------- Front matter ----------------------------
\pagestyle{fancy}
\pagenumbering{roman}

% Cover page, title page and the declaration with the AI-tools box.
\maketitle

\chapter*{Abstract}
\addcontentsline{toc}{chapter}{Abstract}

% Migrated from proposal/proposal.tex (5 October 2026) on 14 October 2026.
Health trackers process data that can reveal sensitive aspects of
a person's life. Research has examined their data practices and users' privacy
concerns, but the way providers describe possible downstream harms needs
closer study. This thesis proposes a content analysis of privacy policies,
store listings and provider websites for 50 widely used apps available in
Germany. Manual walkthroughs will add interface communication.
The analysis will distinguish descriptions of data flows, capabilities such as
inference, and consequences for users. It will examine how harms are presented
and how providers frame privacy protections, guided by literature-based harm
trees.
Coding assisted by a large language model (LLM) will \highlight[orange!15]{be checked against blind
manual coding, with attention to missed mentions. The output will be an
evidence-linked account of harm communication and protection framing in the
captured material, not of data practices, legal compliance or user effects.}\todo{LLM Analysis Methodology is not specified yet}

\noindent\textbf{Keywords:} privacy harms, health trackers, privacy communication, privacy-enhancing technologies, content analysis.



\cleardoublepage
\chapter*{Acknowledgements}
\addcontentsline{toc}{chapter}{Acknowledgements}
\todo[inline]{Do I need this at all?}

\cleardoublepage
\tableofcontents

%---------------------------- Main matter -----------------------------
\cleardoublepage
\pagenumbering{arabic}


% Migrated from proposal/proposal.tex (5 October 2026) on 14 October 2026.
\chapter{Introduction}
\label{ch:introduction}

\section{Problem Statement and Research Gap}
The closest precedent, Malki et al., found inconsistencies between female
health apps' policies and their privacy-related features
\parencite{malki2024exploring}. This thesis instead asks how \emph{health tracking apps} communicate harms across several surfaces, combining harm
coverage, stance and the framing of PETs (Privacy Enhancing Technologies). \highlight[orange!25]{The initial review did not identify any study that combines these elements in this setting. A documented literature search will test and refine this proposed gap.}
\todo{This is a strong claim and should be held until the literature search confirms it with sufficient confidence. The Section should be developped more and more precedent researches to be included.}

\section{Research Questions and Scope}
\noindent\textbf{RQ.} \emph{What types of downstream privacy harms do health
tracker service providers explicitly acknowledge or downplay when describing
their privacy practices or PETs, and how do they frame the protection these
technologies offer?}

\begin{description}[leftmargin=1.5em,itemsep=1pt,topsep=3pt,
                    labelindent=0pt,font=\normalfont\bfseries]
  \item[RQ1] \textit{Coverage.} Which data flows, capabilities and harms are
    mentioned in the captured communication, and which applicable harms
    in the reference model are not found?
  \item[RQ2] \textit{Stance.} How are named harms acknowledged, qualified,
    minimised, denied or presented as the user's responsibility?
  \item[RQ3] \textit{Protection.} What measures and benefits are claimed,
    and how do the scope and strength of these claims relate to the
    reference model?
  \item[RQ4] \textit{Surfaces.} How do coverage, stance and protection claims
    differ between policies, listings, marketing and inspected interfaces?
\end{description}

\noindent\textbf{Scope.} The study covers health tracking apps: apps that let
users record their own health or body data over time, such as activity,
nutrition, sleep, heart rate, mood and reproductive health
(Section~\ref{sec:app-sampling}).
The sample comprises applications from the German storefronts of Google Play and the Apple App Store, as well as the three major applications pre-installed by device manufacturers.
Wearables are considered through their companion apps.

\section{Significance and Expected Contributions}
 \todo{After results/conclusion}

%\noindent\textbf{Users.} A statement about where data is stored may leave the possible
%consequences unclear. Identifying what providers explain can show which
%questions a reader would still need to answer. The US
%post-Roe~\parencite{cao2024deleted} findings motivate attention to fertility data, but do not establish
%the same legal risks for German users.
%
%\smallskip
%\noindent\textbf{Providers and designers.} Comparing surfaces can reveal opportunities
%to explain relevant risks and the limits of protections at decision points.
%
%\smallskip
%\noindent\textbf{Research and policy.}
%\todo[inline]{After Results/Conclusion}

\section{Structure of this Thesis} \todo{After results/conclusion}
%\highlight[orange!25]{Chapter~{\ref{ch:background}} introduces health tracking, privacy harms and
%	the regulatory context. Chapter~{\ref{ch:related-work}} reviews the research
%	that motivates the gap. Chapter~{\ref{ch:methodology}} describes the planned
%	sampling, collection, coding and validation. Chapters~{\ref{ch:results}},
%	{\ref{ch:discussion}} and~{\ref{ch:conclusion}} are reserved for the empirical
%	results, their interpretation and the conclusion. The appendices document
%	the coding scheme, prompts, app sample and worked examples.} 




% Migrated from proposal/proposal.tex (5 October 2026) on 14 October 2026.
\chapter{Background}
\label{ch:background}
\todo{Will be done last.}
%
%In Germany, 73\,\% of smartphone users use health apps to support fitness,
%nutrition or general wellbeing~\parencite{bitkom2025health}. Health trackers can combine information entered by users, such as symptoms
%and cycle dates, with sensor readings and usage records. Together, these data
%may reveal more than users directly provide. A review of 236 studies reports
%that wearable tracker data can support inferences about private attributes,
%including personality traits~\parencite{niksirat2024wearable}.
%
%Data concerning health and a person's sex life are special categories under
%Article~9(1) of the General Data Protection Regulation (GDPR). Activity data
%falls within the health-data definition when
%it reveals information about health status (Article~4(15), Recital~35).
%Articles~5 and~12 require fair and transparent processing and clear information.
%Articles~13 and~14 specify information duties, including purposes, recipients
%and rights~\parencite{gdpr2016}. Guidance allows this information to be layered
%across a policy and interface notices~\parencite{wp292018transparency}. Misleading commercial claims are also subject to
%consumer-protection law~\parencite{ucpd2005}. \highlight[orange!25]{The study does not assume that
%	all communication about harms is voluntary. It asks what providers say about
%	possible consequences alongside their data-practice disclosures.}
%
%A privacy violation and a privacy harm are related but distinct. Citron and
%Solove describe harms such as financial loss, distress and damage to
%relationships~\parencite{citron2021harms}. Knabenhans et al.\ use \emph{harm
%trees} to study the European Digital Identity Framework
%\parencite[\S4.1]{knabenhans2026harms}. A tree links system components, such as
%information disclosed to another party, to capabilities, such as tracking or
%inference, and then to possible harms. Components may arise from the core
%functionality, its digital context or implementation choices. Their analysis
%shows limits of privacy-enhancing technologies (PETs) within that specific
%system and threat model~\parencite{knabenhans2026harms}.
%
%This thesis will adapt the method to trackers. It will distinguish an app's
%intended prediction for its user from an adversary's ability to obtain or use
%that prediction. Martin et al.\ question whether limiting
%third-party access, using inferences instead of raw data, or processing data on
%the device necessarily prevents privacy violations~\parencite{martin2025cookies}.
%The relevant question is which protection is claimed, against which actor and
%under which conditions.



% Migrated from proposal/proposal.tex (5 October 2026) on 14 October 2026.
\chapter{Related Work}
\label{ch:related-work}

\section{Privacy Practices, User Concerns and Risk Communication}
Prior studies examine whether disclosed practices match observed behaviour.
Huckvale et al.\ found that 29 of 36 depression and smoking-cessation apps
transmitted data to Google or Facebook services. Only 12 of 28 apps sending
data to Google and 6 of 12 sending data to Facebook named the respective
company in their policies~\parencite{huckvale2019assessment}. A study of nearly
5{,}000 Android apps found that 81\,\% misrepresented collection or sharing in
the Data Safety section~\parencite{arkalakis2024abandon}. Other work analyses
Apple's privacy labels~\parencite{li2022ios} and fertility trackers through
policy analysis, interface inspection and network measurement
\parencite{shipp2020period,malki2024exploring,hudig2025intimate}.

User studies offer a complementary view. Fitness tracker users considered
harmful data-use scenarios plausible but often unlikely
\parencite{gabriele2020understanding}. US period-tracker users reported privacy
concerns after Roe v.~Wade was overturned~\parencite{cao2024deleted}.
Provider risk communication has also been studied. Zou et al.\ analysed 161
breach notifications drawn at random from those filed with the Maryland
Attorney General between January and June 2018. Hedge
terms such as \emph{may} or \emph{likely} appeared in 112 (70\,\%), and 64
(40\,\%) claimed to have no evidence of misuse, a statement the authors
consider to downplay the risk, because a lack of evidence is not evidence of
absence. Only 9 notifications linked the exposed data to concrete misuse
scenarios, which the authors name as the effective strategy
\parencite[\S5, p.~7]{zou2019might}. Doan et al.\ examined risk and benefit information in
six genetic-testing companies' privacy and research-consent policies
\parencite{doan2024transparency}. Moniz and Almeida gave eight women in
Portugal a design-probe kit on privacy and consent in intimate health
technologies. They conclude that a one-size-fits-all notice leads back to
habituation, and propose that privacy mechanisms should instead inform users
of the specific risks of the data processing, in notices that are tailored to
the context of use, show personalised risk information and use simple language
\parencite[\S5.2]{moniz2025toward}. The authors note that eight participants
may not reflect the target population and that people interested in privacy
are more likely to take part~\parencite[\S3.6]{moniz2025toward}; the proposals
were not tested.

\section{Literature Review}
\todo{To be extended.}



% Migrated from proposal/proposal.tex (5 October 2026) on 14 October 2026.
\chapter{Methodology}
\label{ch:methodology}

%	\section{Draft Status and Study Design}
%	The study plan below is transferred from the proposal dated 5 October 2026.
%	It describes intended work; it does not report a completed study.
%	\todo{Revise this plan against the supervisor meeting of 13 October 2026:
%	remove category quotas; use English primarily with German as fallback;
%	include apps without a privacy policy; extend the keyword space to mental
%	health trackers; exclude connected-device analysis; specify two-coder
%	procedures and justify the sampling and access rules. The ordering and
%	remaining thresholds require a documented decision.}

%%%% Methodology Section Introduction %%%%
This chapter describes how the study will be carried out. The study combines
a structured content analysis~\parencite{krippendorff2019content} with an
adapted cognitive walkthrough~\parencite{wharton1994cognitive}. The sample comprises 50
apps: eligible apps from the German market of Google Play and the App Store,
ranked together by a German-market popularity score based on ratings and reviews, and the health apps that
come preinstalled on the phones of the leading vendor of each platform in Germany. A pilot
with five of these apps will refine the instruments before the main study.
For each app, the privacy policy, the store listings and \highlight[orange!25]{related provider
pages}\todo{Not sure if the providers pages will be collected.} will be collected, and the apps will additionally be inspected in the cognitive walkthroughs.\todo[inline]{Summary of Coding, Analysis, Conclusions, .. after relevant subsections are done} \todo[inline]{Summary of literature review} Figure~\ref{fig:pipeline} summarises this workflow.

\section{Literature Review}
\label{sec:lit-review}
\todo[inline]{Inspired from the paper: Large Language Model Hacking: Quantifying the Hidden Risks of Using LLMs for Text Annotation, cited [5]}
Adding to understanding the related research on the matter, the literature review
supplies the capabilities and harms from which the reference harm trees are
built. It provides also the evidence for the methodological decisions of this
chapter. The review consists of the following the procedure:
fixed queries, the removal of duplicates and of rarely cited works, the
screening of titles and abstracts against explicit criteria, and the extraction
of fixed fields from the included works ~\parencite[App.~A, pp.~56--57]{baumann2025llm}. Afterwards, capabilities are collected to build the harm tree.
\todo[color=green, inline]{done!}

\paragraph{Review questions.} Four questions guide the search, one for each part
of the study. (1)~Which privacy harms and capabilities have been documented for
health trackers, including wearables, starting from
general taxonomies of privacy harm~\parencite{solove2006taxonomy,calo2011boundaries}?
(2)~How have earlier studies sampled apps and analysed their privacy
communication, and with which material, coding scheme and agreement measure,
starting from work on policy analysis~\parencite{wilson2016creation,zimmeck2019maps}?
(3)~How have large language
models been used to annotate privacy policies, and how were they validated?
(4)~How do providers word risk and protection in policies and notices?
\todo{questions may be extended as the thesis evolves}

\paragraph{Stage 0: database search.} The search covers the titles and abstracts
in two open bibliographic databases, OpenAlex~\parencite{priem2022openalex} and
Semantic Scholar~\parencite{kinney2023semantic}.
\todo[inline, color=red]{redo the rest of the literature review, ditching OpenAlex and Sematic Scholar, use traditional platforms.}
% Both index the main venues of
%privacy, security and human--computer interaction research (PoPETs, CCS, USENIX
%Security, IEEE Security and Privacy, CHI and CSCW), the ACM Digital Library, the
%ACL Anthology and arXiv. Baumann et al.\ also searched Scopus, Web of Science and
%Google Scholar. The first two require institutional access to their interfaces,
%and Google Scholar offers no interface for automated queries, so a search in
%them could not be repeated from the project's scripts. Each review question has
%one query, which joins groups of terms with AND and the alternatives within a
%group with OR (Table~\ref{tab:lit-queries}). The domain group, shared by the
%first two queries, names health and fitness apps, fitness trackers, wearables,
%menstrual, fertility and mental-health apps, and femtech. The search period runs
%from 1~January 2014 to the day of the search. The query on language models
%starts on 1~January 2022, as in Baumann et al., because annotation with large
%language models begins then. A script runs the queries and records every hit,
%and the collection is kept unchanged. Works known from earlier reading that meet
%an inclusion criterion but are not retrieved by the queries are added by hand,
%as in the source procedure, and pass through the same filters.
%
%\begin{table}[tbp]
%\centering
%\small
%\caption[Search queries of the literature review]{Search queries of the
%literature review, applied to titles and abstracts in OpenAlex and Semantic
%Scholar. DOMAIN stands for the domain group in the first row.}
%\label{tab:lit-queries}
%\begin{tabularx}{\textwidth}{@{}ll>{\raggedright\arraybackslash}X@{}}
%\toprule
%\textbf{Query} & \textbf{Period} & \textbf{Terms} \\
%\midrule
%\input{tables/literature-queries}% the generated rows end with \bottomrule; a second one here breaks the table
%\end{tabularx}
%\end{table}
%
%\paragraph{Duplicates and citation filter.} Records that share a DOI, an arXiv
%identifier or a normalised title, with publication years at most one year apart,
%are merged. Following Baumann et al., records with fewer than three citations
%per year are then removed. Citations per year are the higher of the two
%databases' counts divided by the number of years since publication, which is at
%least one. The filter keeps the screening manageable and favours works that the
%field has taken up. As in the source procedure, recent works are not exempt, so
%most works from the year of the search are removed; the closest recent studies
%are discussed in the related-work chapter instead.
%
%\paragraph{Stage 1: screening.} Titles and abstracts are screened, and each
%record receives one code. A record is included if it
%\begin{description}
%  \item[I1] documents privacy harms, risks, threats or inferable capabilities
%    that affect the users of health apps, fitness trackers or health wearables;
%  \item[I2] analyses the privacy communication of such apps or wearables, such
%    as privacy policies, privacy labels, store listings, notices or terms;
%  \item[I3] uses or evaluates language models to annotate, classify or extract
%    information from privacy policies or similar privacy text; or
%  \item[I4] analyses how policies, notices or labels present or word data
%    practices, risks or protections, in any domain.
%\end{description}
%A record is excluded if it
%\begin{description}
%  \item[E1] proposes or evaluates a technical system, such as a protocol, an
%    encryption scheme, an architecture or a detection model, without
%    documenting harms or analysing communication;
%  \item[E2] lies outside the domain and topics above;
%  \item[E3] studies health apps or wearables without analysing privacy harms or
%    privacy communication, for example their quality, usability, adoption or
%    features;
%  \item[E4] is not a usable scholarly record, such as a commentary, a
%    dissertation or a duplicate; or
%  \item[E5] concerns language models without annotating privacy text, for
%    example language models as a privacy risk, benchmarks of privacy norms or
%    tools that generate policies.
%\end{description}
%\highlight[orange!25]{Each decision is recorded with its code.}\todo{The screening
%decisions in 01\_Literature/\allowbreak Systematic-Review/\allowbreak screening.csv
%were prepared by an AI assistant from titles and abstracts. Check every decision
%before this section states who screened.}
%
%\paragraph{Stage 2: extraction.} For every included work that is used, a card
%records the review question it answers, the domain, the material, the sampling
%and coding procedure, the agreement statistics, for language-model studies the
%models and their validation, the harms and capabilities named, the findings and
%the authors' stated limits, each with its page, and separates these from the
%reviewer's own critique. A comparison matrix then sets out, for each study, its
%domain, the sources it analyses and how it treats harms; the matrix tests the
%research gap.
%
%\paragraph{Stage 3: capabilities.} Capabilities named in the works included
%under I1 are collected in a table that links each capability to its source, the
%data it needs and the adversary that could use it. The table is the input to the
%harm trees.
%
%\paragraph{Limits.} One person screens, so no inter-rater check of the selection
%is possible. The citation filter removes most works of the last year. The open
%databases lack the abstracts of some records, which are then screened on the
%title alone. The queries are in English, which favours English sources. Later
%literature is added through the change log, with a reason, and does not alter
%the frozen instruments without a logged revision.
%
%\paragraph{Review questions.} Four questions guide the search, one for each part
%of the study. (1)~Which privacy harms and capabilities have been documented for
%health trackers, including wearables, starting from
%general taxonomies of privacy harm~\parencite{solove2006taxonomy,calo2011boundaries}?
%(2)~How have earlier studies sampled apps and analysed their privacy
%communication, and with which material, coding scheme and agreement measure,
%starting from work on policy analysis~\parencite{wilson2016creation,zimmeck2019maps}?
%(3)~How have large language
%models been used to annotate privacy policies, and how were they validated?
%(4)~How do providers word risk and protection in policies and notices?
%
%\paragraph{Sources and search.} Searches cover the proceedings of the main
%venues for privacy, security and human--computer interaction research, namely
%PoPETs, CCS, USENIX Security, IEEE Security and Privacy and CHI/CSCW, together
%with the ACM Digital Library, arXiv and Google Scholar. Queries combine three groups of terms. The first names the
%domain (health app, fitness tracker, wearable, menstrual or fertility app,
%mental-health app). The second names the object (privacy policy, store listing,
%privacy label, data-safety section, breach notification). The third names the
%method or topic (content analysis, cognitive walkthrough, heuristic evaluation,
%harm taxonomy, inference, large language model annotation). German terms are
%used only to find German-market sources. The reference lists and citing papers
%of the closest studies, such as Malki et al.~\parencite{malki2024exploring},
%Hutton et al.~\parencite{hutton2018assessing} and
%Knabenhans et al.~\parencite{knabenhans2026harms}, are searched as well.
%
%
%\paragraph{Selection.} A work is included if it is a peer-reviewed paper, a
%report of an established institution or an archival preprint; is written in
%English or German; and studies mobile health apps or wearables, the analysis of
%privacy communication, harms or risk wording, or the validation of language
%models for coding text. Studies of language models as a privacy risk (for
%example memorisation by chatbots), benchmarks of privacy norms, and tools that
%generate policies are excluded, because they do not inform the coding of
%provider communication. Titles and abstracts are screened first, and the full
%text is read for works that bear on a decision. Each included work is ranked
%\emph{A} (read before the relevant procedure is frozen), \emph{B} (read for a
%specific decision or chapter) or \emph{C} (background). This rank is a reading
%priority and not a quality score. A work whose full text is not freely
%available through open access or the university library is recorded as not
%retrieved, and no claim is based on its abstract alone.
%
%\paragraph{Record and extraction.} The log records for each search round the
%date, the tool, the queries, the exclusions and the works that could not be
%retrieved. Two rounds have been run. On 14 October 2026, a methodology corpus of
%93 works (83 records with a DOI verified against Crossref or DataCite, and 10
%official resources) was compiled from the topics of the supervisor meeting; 89
%are held as full texts. On 23 October 2026, a search for gaps, centred on the
%coding of privacy policies with language models, produced about 75 candidate
%works, of which 46 open-access full texts were obtained. For every work that is
%used, a card records the domain, the material, the sampling and coding
%procedure, the agreement statistics, the findings and the authors' stated
%limits, each with its page, and separates these from the reviewer's own
%critique. A comparison matrix then sets out, for each study, its domain, the
%sources it analyses and how it treats harms; the matrix tests the research gap.
%Capabilities found in the literature are collected in a table that links each
%capability to its source, the data it needs and the adversary that could use it,
%which is the input to the harm trees.
%\todo{Hits and inclusion decisions per query are not yet logged for the venues
%named above; add them to the search log, and update the counts if further rounds
%are run before the instruments are frozen.}
%
%\paragraph{Limits.} One person searches and screens, so no inter-rater check of
%the selection is possible. The search favours English sources and published work
%up to October 2026, and works behind a paywall may be missing. Later literature
%is added through the change log, with a reason, and does not alter the frozen
%instruments without a logged revision.
%
\begin{figure}[t]
\centering
\resizebox{\textwidth}{!}{%
\begin{tikzpicture}[
  font=\footnotesize,
  bx/.style={draw=primary,fill=primary!6,rounded corners=2pt,align=center,
             text width=29mm,minimum height=13mm,inner sep=2.5pt},
  arr/.style={-{Stealth[length=1.6mm]},semithick,accentdark}]
  \node[bx] (a) at (-7.2,0) {\textbf{Pilot and sample}\\store rankings\\preinstalled apps};
  \node[bx] (b) at (-3.6,0) {\textbf{Capture}\\documents for 50 apps\\walkthroughs};
  \node[bx] (c) at (0,0) {\textbf{Code}\\coverage, stance\\protection claims};
  \node[bx] (d) at (3.6,0) {\textbf{Check}\\blind manual coding\\audit missed mentions};
  \node[bx] (e) at (7.2,0) {\textbf{Compare}\\RQ1--RQ4\\with source evidence};
  \draw[arr] (a)--(b); \draw[arr] (b)--(c);
  \draw[arr] (c)--(d); \draw[arr] (d)--(e);
\end{tikzpicture}}
\todo[inline, color=pink]{Draft!}% a margin \todo cannot go inside a float
\caption[Planned workflow]{Planned workflow. The pilot will refine the instruments before they
		are frozen. Later changes will be logged and affected material recoded.}
\Description{Five boxes connected by arrows: pilot and sample, capture,
code, check, compare. Documents cover 50 apps.}
\label{fig:pipeline}
\end{figure}

\section{App Sampling}
\label{sec:app-sampling}
The target sample consists of 50 apps. Previous studies that examined both
app documents and the apps themselves used samples of 20 to 64 apps:
Malki et al.\ installed and inspected 20~\parencite{malki2024exploring},
Huckvale et al.\ analysed 36~\parencite{huckvale2019assessment}, and
Hutton et al.\ evaluated 64~\parencite{hutton2018assessing}.

The sample comprises three groups: \textcolor{accent}{(i)} apps preinstalled
by major device manufacturers; \textcolor{accent}{(ii)} apps listed on
Google Play; and \textcolor{accent}{(iii)} apps listed on the Apple App Store.
Preinstalled apps are selected based on the popularity of each manufacturer's
devices. The remaining places are filled with apps from Google Play and the
Apple App Store using defined collection and ranking rules. No quotas are set for the three source groups \textcolor{accent}{(i)}-\textcolor{accent}{(iii)}. The search frame and eligibility rules were
recorded before collection. The resulting sample composition is reported
in Section~\ref{sec:results-sample}.
\todo[inline, color=green]{done}

\paragraph{Domain.} The study adopts the domain of two published app studies
on health trackers. The first domain is of Hutton et
al.~\parencite{hutton2018assessing}, search categories are derived from the
data types that 105 apps reported most often, because no taxonomy of tracking existed; their health categories are cycling, diet,
exercise, heart, mood, running, sleep, step count and weight. This
survey-based list has no reproductive/fertility category, so the domain of Malki et
al.~\parencite{malki2024exploring} is added, since the paper is about \emph{female health apps}: female reproductive health,
including pregnancy apps. The domain thus covers activity, diet and weight,
heart rate, sleep, mood and reproductive health.

\paragraph{Keywords.} The 15 health keywords of Hutton et al.\ and the 12
terms of Malki et al.\ are searched in the original English and in one German
translation each (Table~\ref{tab:keywords}). A translation names the tracked
data type rather than the word (\emph{eating} becomes \emph{Ernährung}).
Identical spellings and repeated translations are searched once, which gives
49 queries.

% Tables 4.1 and 4.2 share one float on a page of their own ([p]).
\begin{table}[p]
\centering
\begin{singlespace}
% The short form keeps citations out of the List of Tables, where they force
% an extra Biber run.
\caption[Search keywords]{Search keywords. English terms from Hutton et
al.~\parencite[Table~2]{hutton2018assessing} (health categories) and Malki et
al.~\parencite[\S3.1]{malki2024exploring}, each with the German query used.
\emph{workouts} and \emph{menstruation} are spelled alike in both languages and
were searched once.}
\label{tab:keywords}
\small
\setlength{\aboverulesep}{0.25ex}\setlength{\belowrulesep}{0.35ex}
\begin{tabularx}{\textwidth}{@{}>{\raggedright\arraybackslash}p{3.2cm}>{\raggedright\arraybackslash}X>{\raggedright\arraybackslash}X@{}}
\toprule
 & \multicolumn{2}{c}{\textbf{Keywords}} \\
\cmidrule(lr){2-3}
\textbf{Category} & \textbf{English} & \textbf{German} \\
\midrule
Cycling & cycling & Radfahren \\
\addlinespace
Diet & diet & Diät \\
 & eating & Ernährung \\
\addlinespace
Exercise & exercise & Training \\
 & workouts & workouts \\
\addlinespace
Heart rate & heartrate & Herzfrequenz \\
 & heart rate & Puls \\
\addlinespace
Mood & mood & Stimmung \\
 & happiness & Glück \\
\addlinespace
Running & walking & Gehen \\
 & running & Laufen \\
\addlinespace
Sleep & sleep & Schlaf \\
\addlinespace
Steps & steps & Schritte \\
\addlinespace
Weight & weight & Gewicht \\
 & body fat & Körperfett \\
\midrule
Reproductive health & female health, women's health & Frauengesundheit \\
 & reproductive health & reproduktive Gesundheit \\
 & fertility & Fruchtbarkeit \\
 & ovulation & Eisprung \\
 & menstrual, menstruation & menstruation \\
 & period & Periode \\
 & contraception, birth control & Verhütung \\
 & pregnancy & Schwangerschaft \\
 & pregnant & schwanger \\
\bottomrule
\end{tabularx}

\vspace{3\baselineskip}
\caption{Eligibility criteria and how each is checked on the store listing in
the German storefront.}
\label{tab:eligibility}
\setlength{\aboverulesep}{0.4ex}\setlength{\belowrulesep}{0.65ex}
\begin{tabularx}{\textwidth}{@{}l>{\raggedright\arraybackslash}p{3.6cm}>{\raggedright\arraybackslash}X@{}}
\toprule
 & \textbf{Criterion} & \textbf{Check on the listing} \\
\midrule
E1 & Records the user's own health or body data over time, in a category of the
  domain~\parencite{robillard2019availability,malki2024exploring}
  & The description states that the app records, logs or stores the user's own
  measurements or activities, entered by the user or recorded by a sensor or
  paired device. Excluded: instructional and content apps without recording, as
  in Hutton et al.~\parencite{hutton2018assessing}; apps that track a child;
  games, shops and general-purpose apps. A wearable's companion app is eligible
  when the app records or displays the data; the device itself is not analysed. \\
\addlinespace
E2 & Consumer app for the general public
  & The listing offers the app to anyone. Excluded: apps whose listing limits use
  to the members, patients, employees or professional users of an organisation,
  such as a health insurer's service app for its members. \\
\addlinespace
E3 & Free to install
  & App Store: price 0~EUR in the German storefront (Search API field
  \texttt{price}). Google Play: price 0 and marked as free on the product page
  (fields \texttt{price} and \texttt{free}). \\
\addlinespace
E4 & Listing in German or English
  & The listing text in the German storefront is in German or English. \\
\addlinespace
E5 & Not a DiGA, Germany's reimbursable digital health applications
  & A listing that describes prescription by a doctor or reimbursement by
  statutory health insurance is checked against the DiGA directory of the
  Federal Institute for Drugs and Medical Devices (BfArM). \\
\bottomrule
\end{tabularx}
\end{singlespace}
\end{table}

\paragraph{Search.} Each store is searched in a single session, logged out, in
its German storefront: the App Store through Apple's Search
API~\parencite{apple2022itunessearch} and the Google Play store through its
search page, with the collectors in Appendix~\ref{app:scripts}. For each query,
the first 30 results are kept in both stores. A Google Play search page shows
a single page of results, and deeper results need an interface that the store's
\texttt{robots.txt} disallows. Using the same depth for the App Store follows
O'Loughlin et al., who capped both stores at the number of results Google Play
shows so that neither store is favoured~\parencite{oloughlin2019reviewing}. The
App Store is also searched to 200 results, which serves only to check whether
a deeper search would change the sample. On Google Play, the German count used
for the ranking appears only on an app's product page, so the product page of
every Google Play candidate is opened.

\paragraph{Eligibility.} A listing is eligible if it meets the five criteria
of Table~\ref{tab:eligibility}. Each criterion is checked on the store listing
in the German storefront: title, developer, genre, price and description. An
app listed in both stores is decided on one screened listing. If both of its
listings are screened, the app is excluded when either listing fails a
criterion; in the collection, this never happened. A privacy policy is not
required: an app
without one stays in the sample, and the absence is reported. In-app purchases
and subscriptions are recorded but do not exclude an app. Screening follows the
rank order, and each decision is logged with its reason.

\paragraph{Ranking by German-market popularity.} Candidates are ordered by a
\emph{German-market popularity score}, $P_i$, calculated from the feedback
counts of app~$i$ in the German storefronts: ratings count in the App Store and
written reviews count on Google Play. Feedback counts are used because they are the
only count-based measure that is available for the German market. Download and
install counts are either reported worldwide only or not published at all.
%whereas rating and review counts are shown for each storefront.
The score is
therefore a proxy for popularity used to select the sample; it does not measure
users or downloads. Count-based selection has precedent: Sunyaev et
al.\ selected frequently rated health apps and reported a Spearman correlation
of 0.89 between Android rating counts and downloads~\parencite{sunyaev2015availability}.
Their finding supports the use of feedback counts as a proxy.

The available German-market feedback count differs between stores
(Table~\ref{tab:storefront-fields}):
\begin{itemize}
  \item A \textbf{Google Play listing} contributes its German written-review
    count, denoted $v_i$.
  \item An \textbf{App Store listing} contributes its German rating count,
    denoted $r_i$.
\end{itemize}
Thus, a listing supplies one of the two inputs. The unit being ranked is the
\emph{app}, the same app could be present in both stores, or just in one of both. The scoring formula changes accordingly.
 however: its listings are matched across stores before scoring.
An app found only in Google Play has an observed $v_i$ contribution; one found
only in the App Store has an observed $r_i$ contribution. A matched app found
in both stores has both.

A rating and a written review are different feedback actions. They cannot be
treated as interchangeable counts without a conversion. Google Play reviews
serve as the reference unit: one contributes one score unit. Each App Store
rating contributes $c$ units, where the same conversion factor $c$ is applied
to all candidates. The three cases are therefore
\[
  P_i =
  \begin{cases}
    v_i, & \text{Google Play listing only},\\
    c\,r_i, & \text{App Store listing only},\\
    v_i+c\,r_i, & \text{matched listings in both stores}.
  \end{cases}
\]
The three cases can then be written compactly as
$P_i=v_i+c\,r_i$. The addition combines the converted contributions of the two
listings of a matched app; it does not require every app to have both inputs
or imply that an unobserved store count is actually zero.

The conversion is calculated from the matched candidates found in both stores,
before eligibility screening. Let $M$ contain those candidates with positive
counts in both stores. First, calculate the median ratio of Play reviews to
App Store ratings. Then adjust that ratio using the German platform shares:

\[
  m = \operatorname{median}_{j\in M}\!\left(\frac{v_j}{r_j}\right),
  \qquad
  c = m\,\frac{s_{\mathrm{iOS}}}{s_{\mathrm{Android}}}.
\]

Here, $s_{\mathrm{iOS}}$ and $s_{\mathrm{Android}}$ are the platforms' shares of
German mobile web page views in the month preceding collection, taken from
StatCounter~\parencite{statcounter2026os}. These shares approximate the
platform distribution of users. The observed ratio $m$ reflects both the
different platform sizes and the different rates at which users leave ratings
or reviews. Multiplication by the inverse platform-share ratio aims to remove
the platform-size component, so that $c$ converts App Store ratings into Play
review-equivalent units. This assumes that the typical user split of matched
apps follows the platform shares and that platform-specific feedback rates
are sufficiently similar across apps. 

All candidates are ranked by descending $P_i$, with ties broken by the stable
\emph{candidate key} used by the ranking script.

\paragraph{Selection.} The unit of selection is the provider: a listing is skipped when its developer
is already in the sample, through either store or the preinstalled group. In a
trial collection on 22 October 2026\footnote{Collected with the same collection scripts as in Section~\ref{sec:app-sampling} (Search).}, one developer held 10 of the 47 most-rated
Google Play apps under a single privacy policy (Table~\ref{tab:developer-pilot});
RQ1--RQ3 concern what providers
communicate, so further apps of one provider add little new communication and
reduce the number of independent providers. The rule also means that each app
counts once, and both of its listings are captured.

Two listings belong to the same provider if one of three conditions holds.
(P1) They belong to the same developer account of a store, which both stores
identify by a developer ID\@. (P2) Their developer names agree after case,
accents, punctuation, legal forms and regional entity words are removed, as
\enquote{WW International, Inc.} on the App Store and \enquote{WW
INTERNATIONAL, INC} on Google Play do.
(P3) The names differ, but the listings document one provider: the legal entity
that the App Store shows as seller carries the other developer name, or
listings of the two developer accounts link to the same registered web domain,
as seller or developer website, privacy policy or a link in the description.
Daylio, for example, is published by
\enquote{Relaxio s.r.o.} on the App Store and by \enquote{Habitics} on Google
Play; the App Store seller website, daylio.net, is the domain of the Play
privacy policy. P3 matches are made by hand and recorded with their evidence
(Table~\ref{tab:developer-aliases}). A similar name is no evidence:
\enquote{Leap Health} and \enquote{Leapjoy Games} are different developers.
\todo[inline,color=green]{done!}

\paragraph{Preinstalled apps.} Health apps that ship with the phone reach
users without a download, so their rating counts understate their reach and
they cannot be ranked like the store apps. They form a separate group, and
its members follow from a rule with three steps.
\emph{Vendors:} on each platform, the vendor with the largest share of the
platform's mobile page views in Germany is included. The shares come from the
vendor statistics of StatCounter~\parencite{statcounter2026vendor}, read at the
time of the collection and recorded with their date. A vendor's share of its
platform is its share of all mobile page views divided by the platform's share
of all mobile page views~\parencite{statcounter2026os}, so that each vendor is
compared within its own platform. The rule has no threshold and treats iOS and
Android alike, as the ranking does: a German user of either platform most often
meets the configuration of its leading vendor, the same logic as taking the
most widely used apps. Page views approximate the installed base, not sales.
Rules that admit more vendors, such as a minimum share of all page views or
vendors that together cover most of a platform, are reported as checks.
\emph{Apps:} for each leading vendor, the preinstalled apps are identified from
its documentation and the store listings. An app qualifies if the vendor itself
provides it, if it is installed on the vendor's current phones without a
download, and if it meets the eligibility criteria E1 to E5, in particular E1
(it records the user's own health or body data over time). A vendor that ships
no such app contributes none, and this is recorded.
\emph{Treatment:} the store listing of a qualifying app, if it has one, is
removed from the store rankings, and its vendor counts as a provider already in
the sample, so that the vendor's other apps are skipped by the developer rule.
The vendors and apps found are reported in Section~\ref{sec:results-sample}.

\paragraph{Sensitivity checks.} The selection is repeated with one rule changed
at a time: without the developer rule; with equal weights for the two
platforms; on the alternative scale; with the worldwide Google Play rating count
in place of the German review count; at depth 20; with the App Store at depth
200; with each keyword left out; and with English or German keywords only. It
is also compared with an earlier selection that took apps alternately from a
German App Store ranking and a worldwide Google Play ranking. For each variant,
the number of selected providers that change is reported.

Neither store publishes German figures for every listing field, and the two
stores split their fields differently. Apple shows one summary rating per
country or region, tied to the storefront where the customer bought the
app~\parencite{apple2026ratingsoverview}. Google Play has shown phone users
ratings from their registered country since November
2021~\parencite{grinsted2021making}. Download counts are reported only to an
app's developer, in App Store Connect~\parencite{apple2026acquisition} and in the
Play Console~\parencite{google2026playstatistics}. These sources do not say
which public field follows which rule, so Table~\ref{tab:storefront-fields}
rests on a field check. On 22 October 2026 it requested the same listings from
the German, US and British storefronts and compared the values from Apple's
search and lookup endpoint~\parencite{apple2022itunessearch}, the App Store
product pages, and the Google Play product pages read with an open-source
parser~\parencite{jomingyu2024googleplayscraper}. On the App Store, every rating
figure describes the requested storefront and no worldwide total is published.
On Google Play, the rating count and the install figures were the same in every
storefront, while the average rating and the one-to-five-star breakdown
changed: the breakdown is the country's distribution scaled to the worldwide
count. Both stores show an app's position in a category chart of the
storefront. The check covered few apps in one session, so it shows how the
fields behave and does not enter the sample. The two stores therefore count
ratings on different bases.

\begin{table}[t]
\centering
\begin{singlespace}
\caption[Public listing data by geographic basis]{Public listing data by geographic basis. Observed on 22 October 2026
for up to six App Store and five Google Play listings in the German (DE), US
and British (GB) storefronts; fields shown only on product pages were checked
for one app per store in DE and US\@. \emph{Market-specific}: the value changed
with the storefront. \emph{Worldwide}: it was identical across storefronts.
}
\label{tab:storefront-fields}
\small
\setlength{\aboverulesep}{0.25ex}\setlength{\belowrulesep}{0.35ex}
\begin{tabularx}{\textwidth}{@{}>{\raggedright\arraybackslash}p{2.3cm}>{\raggedright\arraybackslash}X>{\raggedright\arraybackslash}X@{}}
\toprule
\textbf{Field} & \textbf{App Store} & \textbf{Google Play} \\
\midrule
Rating count & \textbf{Market-specific}, also for the current version; no
  worldwide total. Yazio: 446{,}768 (DE), 51{,}379 (US), 12{,}351 (GB)
  & \textbf{Worldwide}, labelled \enquote{reviews} on the page; no country
  count. Home Workout: 4{,}045{,}258 (DE, US), 4{,}045{,}260 (GB) \\
\midrule
Average rating & \textbf{Market-specific}, also for the current version.
  Yazio: 4.62 (DE), 4.70 (US)
  & \textbf{Market-specific}. MyFitnessPal: 3.94 (DE), 4.41 (US) \\
\midrule
One-to-five-star breakdown & \textbf{Market-specific}, on the product page
  only. Yazio five-star: 349{,}509 (DE), 42{,}759 (US)
  & \textbf{Mixed}: the country's distribution scaled to the worldwide count.
  MyFitnessPal one-star: 428{,}497 (DE), 205{,}443 (US) \\
\midrule
Written reviews & A few review texts on the product page; \textbf{no public
  count}
  & Undocumented embedded count, not displayed, \textbf{market-specific}: it
  follows the storefront country, not the interface language (25 October).
  MyFitnessPal: 32{,}560 (DE), 454{,}136 (US) \\
\midrule
Downloads or installs & \textbf{Not public}; reported to the developer
  only~\parencite{apple2026acquisition}
  & \textbf{Worldwide} band on the page (\enquote{100M+}) and an
  undocumented exact count (Home Workout: 303{,}142{,}208 in all three);
  country figures to the developer only~\parencite{google2026playstatistics} \\
\midrule
Chart position & \textbf{Market-specific} badge on the product page (Yazio:
  no.~3 in Health \& Fitness, DE) and public category chart pages
  & \textbf{Market-specific} badge when ranked (MyFitnessPal: no.~2 top
  grossing in Health \& Fitness, US; none in DE) \\
\midrule
Search result order & Endpoint order per storefront; not verified as the
  native store rank
  & Page order per storefront; not verified as the native store rank \\
\bottomrule
\end{tabularx}
\end{singlespace}
\end{table}

\section{Material and Language} For each app, the core corpus will include
the policy, store descriptions, structured store disclosures, the main product
page, and privacy/security pages linked directly from it or the policy. A
fixed supplementary search will capture provider-hosted privacy blog posts
from the preceding 12 months. Ads encountered during collection will form an
opportunistic supplement. The absence of ads or blog posts will not count as
absence of a harm in the core corpus.

Language is recorded per source. English text will be primary for every
captured source, including interface text, and German will be used when no
English version is available; results describe this English-first corpus. If
both versions exist, both will be captured and their differences logged.
Original wording will be coded. The rule does not apply to the search:
keywords combine German and English terms, and the storefronts remain German.
Each source will be archived with URL, date, language, file hash and app
version where applicable. Shared policies will be coded
once and linked to all relevant apps.

\section{App Walkthroughs} Twelve distinct apps will be inspected, six per domain
with three Android and three iOS apps in each. Selection will follow the
ranking, alternating Android and iOS slots without reusing an app; the balance
serves coverage, not platform comparison. Apps requiring an unavailable paired
device, payment or identity verification will be skipped with a recorded
reason, and unfilled slots reported. Each session will use a fresh install,
synthetic data and a study account where required.

Tasks will cover finding collected data, inferred data, breach consequences
and incident responses, reaching the policy, stopping optional third-party
sharing, exporting data and deleting the account. At each step the researcher
will ask Wharton et al.'s four questions: would the user pursue the intended
goal, notice the action, connect it to that goal and recognise progress
\parencite{wharton1994cognitive}? The pilot will fix starting points and
time limits. Outcomes will distinguish completion, information not found,
unsupported tasks and access barriers. Screen recordings will preserve the
text encountered. This is an expert inspection, not a test of real users.

\section{Reference Harm Trees} Before main coding, a tree will be built for
each domain from prior tracker research. For fitness, sources are a survey of
tracker inferences, threats and vulnerabilities~\parencite{niksirat2024wearable}
and users' harm scenarios~\parencite{gabriele2020understanding}. For fertility,
sources are analyses of policies, data flows and interfaces
\parencite{shipp2020period,malki2024exploring,hudig2025intimate} and post-Roe
user concerns~\parencite{cao2024deleted}. Sources found in the logged search
will be added with a recorded reason. Each pathway will name the data, potential
adversary, access conditions, capability and possible consequence. Required
conditions will be connected explicitly, following Knabenhans et al.'s
AND/OR approach. A prediction made for the user will not itself establish
third-party access (Appendix~\ref{app:example}). As in Knabenhans et al., capabilities include incapabilities, such as a user
being unable to use a feature without agreeing to data sharing.

The component origins will distinguish core functionality, digital context
and implementation choices within the stated model. For every app and
pathway, an applicability record will note the required feature, its
evidence (listing, policy or walkthrough), the access and context
assumptions, the decision (applicable, not applicable or uncertain) and the
reason. A harm category applies to an app when at least one of its pathways
does. The trees are reference models of possible harms, not complete
descriptions of the apps' actual systems. Uncertain cases will be counted and
reported, not silently dropped, so that sparsely documented providers do not
vanish from the gap analysis.

\section{Harm Categories} The starting categories are Knabenhans et al.'s
physical, economic, psychological, autonomy/\allowbreak deprivation, discrimination,
reputational and societal harms~\parencite[App.~A]{knabenhans2026harms},
drawn from Citron and Solove and the CyberPeace Institute. Following them, a
harm is an impairment of the ability to function as one otherwise would in
one's usual context~\parencite[\S4.1]{knabenhans2026harms}; here, the baseline
is the user's situation without the modelled pathway. A relationship tag will
retain Citron and Solove's relationship harms~\parencite{citron2021harms},
which the list omits. New harms found in the pilot will be documented rather
than forced into a category.


\section{Pilot} The pilot takes place after sampling and after the reference
harm trees have been built. Five apps, three fitness and two fertility apps,
will be drawn at random from the sample with a recorded seed. The pilot will
test the boundaries of collection and coding and estimate capture time, coding
units, manual effort and model cost. It will also settle the choices left open
in the previous sections: the starting points and time limits of the
walkthroughs, the stance definitions, the handling of new harms, and the LLM
and its prompts. The instruments will then be frozen. The pilot apps will be
recoded with the final instruments and remain in the sample, but will not
enter the validation set.

\section{Coding Units and Axes} To code a text is to attach labels that
record what it says. A unit will be one paragraph or one coherent block of
interface text. Adjacent text will be available for context, but supporting quotations must
be identified precisely. The codebook will define each label, a fitting
example and a counter-example. Table~\ref{tab:axes} summarises the three axes.
Appendix~\ref{app:example} shows the procedure on a constructed example.

\begin{table}[t]
\centering
\begin{singlespace}
\caption[Coding axes and boundaries]{Coding axes and boundaries. All example wording is constructed.
Several labels may apply to one unit.}
\label{tab:axes}
\small
\setlength{\aboverulesep}{0.25ex}\setlength{\belowrulesep}{0.35ex}
\begin{tabularx}{\textwidth}{@{}L{1.6cm}L{2.8cm}>{\raggedright\arraybackslash}X>{\raggedright\arraybackslash}X@{}}
\toprule
\textbf{Axis} & \textbf{Label} & \textbf{What it records} & \textbf{Example or boundary} \\
\midrule
A: Coverage & Component & Data collected, stored or disclosed, or a system requirement
  & \enquote{We store your cycle dates.} \\
\rowrule
 & Capability & Distinguishing users, tracking/profiling, inferring an attribute, or an incapability
  & \enquote{We estimate your fertile window.} \\
\rowrule
 & Concrete harm & A named negative consequence and its category
  & \enquote{Disclosure could cause distress.} \\
\rowrule
 & Generic risk & Possible misuse without a concrete consequence
  & \enquote{Your data could be used against you.} \\
\rowrule
 & No relevant content & None of the coverage labels applies
  & A contact address alone \\
\rowrule
 & Additional details & Breach consequences, incident response and relationship impacts
  & Code concrete response steps separately from a general legal-notice promise \\
\midrule
B: Stance & Acknowledged & A named harm presented as possible
  & \enquote{A disclosure could cause distress.} \\
\rowrule
 & Qualified & Uncertainty or likelihood language, with its target recorded
  & \enquote{A leak is unlikely, but could cause distress.} \\
\rowrule
 & Minimised & A named harm dismissed without supporting reasons in the unit or adjacent text
  & \enquote{Such distress is no cause for concern.} \\
\rowrule
 & Responsibility shifted & Responsibility for a named harm placed on the user
  & \enquote{You alone must prevent harassment caused by sharing your location.} \\
\rowrule
 & Denied/prevented & A named harm presented as impossible or prevented
  & \enquote{This makes harassment impossible.} \\
\midrule
C: Protection & Measure and scope & Claimed measure, actor, data, outcome and conditions
  & On-device storage is a measure. A no-sale pledge is a practice commitment. \\
\rowrule
 & Framing and strength & Data-flow or harm framing, with partial, scoped or absolute protection
  & \enquote{This reduces exposure} differs from \enquote{No one can harm you.} \\
\rowrule
 & Reassurance & Broad safety language without a specific measure
  & \enquote{Your data is safe with us.} \\
\rowrule
 & Potential mismatch & Human-reviewed tension between the claim and a pathway within its scope
  & A residual risk outside the claim's scope is not sufficient \\
\bottomrule
\end{tabularx}
\end{singlespace}
\end{table}

\subsection{Axis A: Coverage} All explicitly mentioned nodes will be
recorded, including whether each is asserted, presented as possible or denied.
The highest level reached will be a separate summary, so a harm
mention does not erase accompanying data-flow information. Feature descriptions can mention inference without presenting it as a risk.

\subsection{Axis B: Stance} Stance labels will be linked to the
particular harm and supporting span, rather than assigned to the paragraph
as a whole. A statement may acknowledge a harm and qualify its likelihood,
and the coder will distinguish uncertainty about a breach from uncertainty
about its consequences. Minimisation will be judged from wording in the unit
or adjacent text, never from inferred intent; qualification alone does not
establish it. RQ2 thus covers stance toward named harms; dismissive language
without one is coded as reassurance. These distinctions adapt Zou et
al.~\parencite{zou2019might} and will be refined in the pilot.

\subsection{Axis C: Protection} The coder will record the measure
and claimed benefit, with actor, data and conditions where specified. Any
technical measure claimed to limit exposure or misuse counts, PET-labelled or
not, including ordinary security controls; non-technical commitments are kept
separate. A \emph{potential mismatch} requires a pathway within the claim's
scope and an explanation of why the measure misses its necessary conditions;
partial protection is not rejected merely because other risks remain. Vague
scope will be recorded as unclear, and no deployed system judged safe or
unsafe.

\section{LLM-Assisted Coding} Polisis and recent LLM approaches are precedents for
data-practice analysis~\parencite{harkous2018polisis,rodriguez2024llm}. The
harm-focused scheme will use prompts and a model selected on the pilot, then
frozen~\parencite{tornberg2024best}; an open-weight model will be preferred if
it performs close to a hosted one, keeping the run repeatable. Prompts will
work in stages: relevant spans, then stance for detected harms, then
protection details. All retained harm mentions and potential mismatches will
be reviewed manually, so the model acts as a screener whose prompts favour
recall. The author will decide tree applicability and mismatches, logging
evidence and reasoning; uncertain outputs or missing spans enter manual
review rather than becoming negative codes. Each unit will run three times;
agreement is a stability check, not validation, and labels returned once are
reviewed rather than discarded. Raw output is kept apart from adjudicated
labels, so performance is not measured after correction.

\section{Validation and Missed Mentions} Ten apps, five per domain, will be
drawn at random with a recorded seed from providers and documents not used in
the pilot, stratified by the language of the coded policy where both occur;
the numbers of independent providers and unique documents will be reported.
The author will code them fully before seeing model output, and they will not
be used to tune prompts. Because prompt wording alone can change
conclusions~\parencite{baumann2025llm}, two paraphrased prompts will also be
run and headline counts compared. Disagreements will be adjudicated with each
label's source hidden; two apps will be recoded after at least two weeks,
before drafting, to estimate intra-coder consistency. A second coder will be
sought for applicability, minimisation and mismatch decisions; failing that,
the supervisor will review a sample, including judged non-mismatches, with
disagreements recorded before discussion. A detection is correct only with
the right label, unit and evidence span, and the staged pipeline is scored end
to end, so first-stage misses count. Reported will be per-code precision
(share of model labels supported), recall (share of manual mentions found)
and model--human Cohen's $\kappa$~\parencite{cohen1960coefficient}, with
counts, prevalence and Gwet's AC1, since rare codes can depress $\kappa$
despite high agreement~\parencite{feinstein1990high,gwet2008computing}.
App-level detection is also reported, by domain and, where counts allow, by
source and language. Exploratory 95\,\% intervals will resample providers to
respect shared sources; with ten apps they will be wide.

A separate missed-mention audit, reported apart from validation, will search
the remaining corpus for harm-related vocabulary in both languages and check
every unlabelled hit manually. To expose the vocabulary's blind spots, 100
units outside the hits (or all, if fewer), allocated across domains, languages
and sources in advance, will also be read. Recall will be reported as
unestimable where manual positives are too few. For any code with fewer than
20 validation positives or recall below 0.90, all relevant material will be
checked manually or corpus-wide absence claims withheld; these are review
triggers, not proof of accuracy.

\section{Analysis and Reporting} Results will be descriptive. For RQ1, each harm's denominator
will include only applicable apps with the required core sources captured. A
source the provider does not offer is recorded as such; only sources that
exist but cannot be retrieved reduce the denominator. Inapplicable pathways
and text with no identified mention will have different labels. Counts will be per app and surface, with
shared-provider results also shown after grouping repeated providers.
RQ2 will report harm--stance combinations and their denominators. RQ3 will
summarise measures, scope, claim strength and manually reviewed mismatches.
Examples that support and challenge the expected gap will both be retained.

For RQ4, the principal policy--interface comparison will use the same 12
inspected apps; the 50-app document analysis gives the broader overview, and
other surface pairs report how many apps have both surfaces. At the fixed
decision points, onboarding before consent and the privacy-settings screen,
text shown directly, behind a visible link or not found within the task limit
will be distinguished. Inconsistent claims will be reported apart from mere
differences in detail. Absence will mean \enquote{no mention was identified
in the captured material}, with the instrument's measured limitations stated.

\section{Ethics and Feasibility} No human participants are planned.
Walkthroughs use synthetic data, store terms will be checked before automated
retrieval, and any review requirements confirmed before collection. The
pilot's estimates will fix a minimum feasible study (core documents, the blind
validation set, the walkthroughs) before collection; if time runs short,
supplementary marketing collection is reduced first and other changes are
agreed with the supervisor.

\section{Limitations} The sample favours free apps with high German-market
popularity scores found by the chosen keywords and does not represent all
health-tracking apps or their users in Germany; one app per provider also
leaves out further apps of large providers. Feedback counts accumulate over
different app lifetimes and depend on prompts to rate or review. The common
scale assumes a typical platform distribution and approximately stable
platform-specific feedback rates across apps. StatCounter's page-view shares
are an approximation of the user distribution, and a listing not found in one
store contributes no count from that store. These limits prevent interpreting
the score as a count of users or communication encounters. The corpus is a
dated snapshot, walkthroughs cover selected paths, and reference trees depend
on the chosen literature. Unless a second coder joins, one person provides the
reference labels, so model agreement is not independent human-coder
reliability~\parencite{mcdonald2019reliability}. Missed mentions may remain.
The study cannot establish provider intent, actual harms or effects on users.



\chapter{Results}
\label{ch:results}

% Report findings only. Interpretation belongs in the discussion.

\section{Literature Search}
\label{sec:results-lit}
This section reports the literature search of Section~\ref{sec:lit-review}.
The counts come from the search script and the screening file; the included
works are listed in Appendix~\ref{app:literature}.
\todo{The screening decisions in 01\_Literature/\allowbreak Systematic-Review/\allowbreak screening.csv
were prepared by an AI assistant from titles and abstracts. Check every
decision before this section states who screened.}

\begin{table}[tbp]
\centering
\small
\caption[Records at each stage of the literature search]{Records at each stage
of the literature search of 26 October 2026. Hits are counted per query and
database before duplicates are removed.}
\label{tab:lit-flow}
\begin{tabular}{@{}lrr@{}}
\toprule
 & \textbf{OpenAlex} & \textbf{Semantic Scholar} \\
\midrule

% Generated by scripts/literature-search/v001/report.py from the search of 2026-10-26 and its screening; do not edit by hand.
Q1 & 1{,}997 & 1{,}355 \\
Q2 & 991 & 731 \\
Q3 & 117 & 76 \\
Q4 & 249 & 151 \\
Manually added & \multicolumn{2}{c}{15} \\
\midrule
Hits & \multicolumn{2}{c}{5{,}682} \\
Records after removing duplicates & \multicolumn{2}{c}{3{,}506} \\
Removed by the citation filter & \multicolumn{2}{c}{2{,}553} \\
Records screened (titles and abstracts) & \multicolumn{2}{c}{953} \\
Excluded, E1 & \multicolumn{2}{c}{248} \\
Excluded, E2 & \multicolumn{2}{c}{236} \\
Excluded, E3 & \multicolumn{2}{c}{283} \\
Excluded, E4 & \multicolumn{2}{c}{6} \\
Excluded, E5 & \multicolumn{2}{c}{12} \\
Included, I1 & \multicolumn{2}{c}{76} \\
Included, I2 & \multicolumn{2}{c}{46} \\
Included, I3 & \multicolumn{2}{c}{13} \\
Included, I4 & \multicolumn{2}{c}{33} \\
Included in total & \multicolumn{2}{c}{168} \\
\bottomrule

% the generated rows end with \bottomrule; a second one here breaks the table
\end{tabular}
\end{table}

\paragraph{Search.} OpenAlex and Semantic Scholar were searched on 26 October
2026 at 18:13~UTC\@. The four queries returned 5{,}667 hits, 3{,}352 of them for
the query on harms (Table~\ref{tab:lit-flow}). Fifteen works were added by
hand. Merging duplicates left 3{,}506 records.

\paragraph{Citation filter.} The filter removed 2{,}553 records, among them 611
of the 648 records published in 2026, and left 953 records for screening. Ten
of the fifteen works added by hand passed it. For 111 of the 953 records
neither database held an abstract, so they were screened on the title.

\paragraph{Screening.} Screening included 168 records: 76 that document harms
or capabilities (I1), 46 that analyse the privacy communication of health apps
or wearables (I2), 13 on language models and privacy text (I3) and 33 on the
presentation and wording of policies, notices or labels (I4). Of the 785
excluded records, 248 proposed technical systems (E1), 236 lay outside the
domain (E2) and 283 studied health apps or wearables without analysing privacy
(E3); 6 were not usable records (E4) and 12 concerned language models without
annotating privacy text (E5). The included works were published between 2014
and 2026: 55 in 2014--2019, 45 in 2020--2022 and 68 in 2023--2026. Eight have no
DOI and are linked to their database page.

\section{Sample}
\label{sec:results-sample}
This section reports how the sample of Section~\ref{sec:app-sampling} was
drawn. All figures come from the collector outputs and the ranking script
(Appendix~\ref{app:scripts}); the sample is listed in
Appendix~\ref{app:app-sample}.
\todo{The screening decisions in 04\_Study/\allowbreak App-Sample/\allowbreak screening.csv were
prepared by an AI assistant from the store listings. Check every decision
before this section states who screened.}

\paragraph{Collection.} The App Store was searched on 23 October 2026 between
19:46 and 19:53~UTC\@. Of the 49 queries, 48 returned between 64 and 190 results;
\emph{reproduktive Gesundheit} returned none. The first 30 positions held 1{,}440
hits and 848 distinct listings. Google Play was searched on 24 October 2026
between 15:07 and 15:29~UTC\@. Each query returned between 8 and 50 results, 1{,}178
in total; the first 30 positions held 1{,}138 hits and 645 distinct apps, and
every hit carried an app ID\@. The product pages of all 645 apps were read on
24 and 25 October 2026; one page that timed out was read again, so none is
missing.

\paragraph{Preinstalled apps.} The vendor shares for September 2026, read on
22 October 2026, give Samsung 39.3\,\%, Apple 27.6\,\% and Xiaomi 21.4\,\% of
German mobile page views, and Huawei and Google about 1.5\,\% each
\parencite{statcounter2026vendor}. Apple is the only vendor of iOS (27.7\,\% of
page views). On Android (72.3\,\%), Samsung leads with 54\,\% of the
platform's page views, followed by Xiaomi with 30\,\%
\parencite{statcounter2026os}. The leading vendors are therefore Apple and
Samsung. Their qualifying apps are Apple Health and Apple Fitness on the iPhone
and Samsung Health on Samsung phones. The rule therefore yields three
preinstalled apps, which leaves 47 places for the ranking. A threshold of 10\,\%
of all page views would admit Xiaomi as well, but no preinstalled tracker could
be identified on Xiaomi's international software, so the three apps are the same
under that rule and under a rule that covers 80\,\% of Android. All three have store listings, which
were removed from the ranking; Apple Health, for instance, has 2{,}134 ratings
in the German App Store.

\paragraph{German-market popularity score.} Of the candidates, 125 apps were
found in both stores. The calibration used the 116 with positive rating and
review counts in both stores. Their counts correlated at Spearman 0.79, and
the median ratio was $m = 0.233815$. With the German platform shares of
September 2026, Android 72.3\,\% and iOS
27.7\,\%~\parencite{statcounter2026os}, the conversion factor was
\[
  c = 0.233815\times\frac{0.277}{0.723}\approx 0.089581.
\]
Thus, one App Store rating contributed approximately 0.09 score units and one
Play review contributed one unit. For example, komoot had 417{,}443 German App
Store ratings and 99{,}884 German Play reviews. Using the unrounded factor,
its score was $P_{\mathrm{komoot}}\approx137{,}279$. This is a combined feedback
score, not a user count; the script ranked candidates with unrounded scores.
Across all 644 Google Play candidates, the German review count and the
worldwide rating count correlated at 0.95, yet they rank the leading apps
differently: among the 100 apps with the most installs worldwide, the German
review count ranged from 5 to 99{,}884.

\paragraph{Screening and selection.} Table~\ref{tab:sample-flow} gives the
cascade. Merging the two stores' listings gave 1{,}365 candidates. In descending
order of German-market popularity score, 70 candidates were reached before
the 47 places were filled.
Of these, 11 were excluded at screening because the app records no health or
body data over time (E1), and 12 were skipped because their developer was
already in the sample. Developer names that differ between the two stores were
matched by hand for 13 providers, among them Clue, Pregnancy+, Fddb and
Headspace, each on the evidence listed in Table~\ref{tab:developer-aliases};
no store account was found under two providers, and no two providers among the
70 candidates reached named the same web domain. Three apps whose titles were
translated between the stores were paired by hand. Across the main selection and all sensitivity checks, 275
listings were screened: 156 were eligible and 119 excluded (E1~114, E2~1,
E3~2, E4~1, E5~1).

\begin{table}[t]
\centering
\begin{singlespace}
\caption[From search hits to the 47 ranked apps]{From search hits to the 47 ranked apps. Positions are the first 30 per
query; the same app in both stores is one candidate.}
\label{tab:sample-flow}
\small
\begin{tabular}{@{}lrr@{}}
\toprule
 & \textbf{App Store} & \textbf{Google Play} \\
\midrule
Queries & 49 & 49 \\
Hits within the first 30 positions & 1{,}440 & 1{,}138 \\
Distinct listings & 848 & 645 \\
Listings of preinstalled apps removed & 2 & 1 \\
Listings ranked & 846 & 644 \\
\midrule
\multicolumn{3}{@{}l}{\textbf{Candidates after merging the stores}} \\
\quad in both stores & \multicolumn{2}{r}{125} \\
\quad App Store only & \multicolumn{2}{r}{721} \\
\quad Google Play only & \multicolumn{2}{r}{519} \\
Reached in descending score order & \multicolumn{2}{r}{70} \\
\quad excluded at screening (E1) & \multicolumn{2}{r}{11} \\
\quad skipped: developer already in the sample & \multicolumn{2}{r}{12} \\
\quad selected & \multicolumn{2}{r}{47} \\
\bottomrule
\end{tabular}
\end{singlespace}
\end{table}

\paragraph{The sample.} The 47 ranked apps come from 47 providers; 33 are listed
in both stores, 8 only on Google Play and 6 only on the App Store. With Apple
Health, Apple Fitness and Samsung Health, the sample holds 50 apps. By primary
category, the ranked apps comprise 22 activity apps, 12 diet and weight apps, 4
cycle trackers, 2 pregnancy apps, 3 sleep apps, 2 mood apps, 1 heart-rate app and
1 wearable companion. The selected apps have between 648 and 447{,}096 German App
Store ratings (39 apps with an App Store listing) and between 1{,}486 and
99{,}884 German Google Play reviews (41 apps with a Google Play listing). All are
free to install, and all 41 Google Play listings name a privacy policy. Of the
47 apps, 42 were found within the first ten positions of at least one query and
5 within positions 11 to 20.

\paragraph{Sensitivity.} Table~\ref{tab:sensitivity} shows how many of the 47
selected providers change when one rule is changed at a time. Without the
developer rule, Leap Fitness Group would take 7 of the 47 places, Simple Design
3, and Google and Nike 2 each (Table~\ref{tab:developer-rule}).

\begin{table}[t]
\centering
\begin{singlespace}
\caption[Sensitivity of the selection]{Sensitivity of the selection. Number of the 47 selected providers that
are not selected when one rule changes.}
\label{tab:sensitivity}
\small
\begin{tabularx}{\textwidth}{@{}Xr@{}}
\toprule
\textbf{Variant} & \textbf{Providers changed} \\
\midrule
No developer rule (the 47 places go to 37 providers) & 10 \\
Equal weights for iOS and Android & 5 \\
Alternative scale (share of each store's total count) & 0 \\
Worldwide Google Play rating count instead of German reviews & 11 \\
Earlier selection: store alternation with the worldwide Google Play count & 10 \\
Depth 20 instead of 30 & 1 \\
App Store searched to depth 200 & 2 \\
English keywords only & 2 \\
German keywords only & 3 \\
One keyword left out (49 runs; largest change) & 1 \\
\bottomrule
\end{tabularx}
\end{singlespace}
\end{table}

\section{What Privacy Communication Covers}
\todo[inline]{To be written.}

\section{What Privacy Communication Omits}
\todo[inline]{Inference, breaches, incident handling.}

\section{Differences Across Platforms and Categories}
\todo[inline]{To be written.}

\section{Privacy Content in Advertising}
\todo[inline]{To be written.}



\chapter{Discussion}
\label{ch:discussion}

\section{Answering the Research Questions}
\todo[inline]{To be written.}

\section{Implications for Users}
\todo[inline]{To be written.}

\section{Implications for Providers and Designers}
\todo[inline]{To be written.}

\section{Implications for Regulation}
\todo[inline]{To be written.}

\section{Threats to Validity}
\todo[inline]{To be written.}



\chapter{Conclusion}
\label{ch:conclusion}

\section{Summary}
\todo[inline]{To be written.}

\section{Future Work}
\todo[inline]{To be written.}



%---------------------------- Back matter -----------------------------
\cleardoublepage
\pagestyle{fancy-lists}

% Unnumbered lists chapter, as in the template: \phantomsection lets hyperref
% point to the right page, \sectionmark sets the running head.
\phantomsection
\addcontentsline{toc}{chapter}{\listoftitlename}%
\chaptermark{\listoftitlename}

\phantomsection
\addcontentsline{toc}{section}{\listfigurename}%
\sectionmark{\listfigurename}
\listoffigures
\cleardoublepage

\phantomsection
\addcontentsline{toc}{section}{\listtablename}%
\sectionmark{\listtablename}
\listoftables
\cleardoublepage

\phantomsection
\addcontentsline{toc}{section}{\lstlistlistingname}%
\sectionmark{\lstlistlistingname}
\lstlistoflistings
\cleardoublepage

\phantomsection
\addcontentsline{toc}{section}{\bibname}%
\sectionmark{\bibname}
\printbibliography[title={\bibname}]
\cleardoublepage

\pagestyle{fancy}
\appendix

\chapter{Codebook}
\label{app:codebook}
\todo{Full coding scheme with definitions and examples, from codebook v3.}

%----------------------------------------------------------------------
\section{Worked Example}
\label{app:worked-example}

This section walks through the whole procedure once, on one privacy policy,
using codebook version~3. The app, \enquote{CycleMate}, and its policy are
constructed: the text is written so that every important coding decision
occurs at least once. Real policies are expected to stay mostly at the level
of data practices and to name few harms.

\subsection*{How to read this example}

A few terms are used throughout. They are explained here in plain words.

\begin{description}[leftmargin=0pt,itemsep=2pt]
  \item[Unit.] One paragraph of the policy. Coding happens unit by unit.
  \item[Code.] A short label that records what a unit \emph{says}, never what
    the provider might mean. Every code needs a quotation from the unit
    (its \emph{evidence span}).
  \item[Three questions (axes).] \textbf{A}: What is mentioned? \textbf{B}: If
    a harm is named, how is it presented? \textbf{C}: What protection is
    promised?
  \item[Reference tree.] A list of the ways harm \emph{could} happen with a
    fertility tracker, written down before any policy is read. Each way is a
    \emph{pathway}: data and access conditions, then something an actor
    can do with them (a \emph{capability}), then the harm.
  \item[Applicable.] A pathway is applicable to an app if the app has the
    features the pathway needs. Only applicable harms can be missing from the
    app's communication.
  \item[Potential mismatch.] A protection claim that, read literally, covers
    a pathway which the named measure does not actually block. Only the author
    decides this, never the language model.
\end{description}

\subsection*{Step 0: The reference tree}

The tree is built once for the fertility domain, from the tracker
literature, following the harm-tree method of Knabenhans et
al.~\parencite{knabenhans2026harms}. \Cref{tab:we-tree} shows five of its
pathways. In the \emph{Conditions} column, AND means that every condition is
needed; OR means that any one of them is enough. Each condition carries its
origin: \emph{core} (the app cannot work without it), \emph{context} (it comes
from running on a phone and network), or \emph{implementation} (the provider
chose it and could choose otherwise).

\begin{table}[H]
\centering
\caption{Excerpt of the reference tree for fertility trackers.}
\label{tab:we-tree}
\small
\begin{tabularx}{\textwidth}{@{}l>{\raggedright\arraybackslash}X
                                   >{\raggedright\arraybackslash}p{3.1cm}
                                   >{\raggedright\arraybackslash}p{3.3cm}@{}}
\toprule
\textbf{ID} & \textbf{Conditions (origin)} & \textbf{Capability} & \textbf{Harm} \\
\midrule
F1 & Logged cycles (core) AND another person can open the unlocked phone
     (context)
   & Partner infers a possible pregnancy (\texttt{K-INF})
   & Unwanted disclosure causes distress (\texttt{H-PSY}); relationship tag
     \texttt{T-REL} \\
\addlinespace
F2 & Logged cycles (core) AND [cloud backup copy (implementation) OR seized
     phone (context)]
   & Authorities infer a pregnancy or pregnancy loss (\texttt{K-INF})
   & Legal pressure, restricted reproductive choices (\texttt{H-AUT}) \\
\addlinespace
F3 & Device identifier (context) AND analytics or advertising SDK
     (implementation)
   & Profiling across apps (\texttt{K-TRK}); user sorted as
     \enquote{trying to conceive} (\texttt{K-DIS})
   & Different prices or offers (\texttt{H-DSC}) \\
\addlinespace
F4 & Predictions offered only if the user agrees to data sharing
     (implementation)
   & User cannot use the core feature without sharing (\texttt{K-INC})
   & Restricted choice (\texttt{H-AUT}) \\
\addlinespace
F5 & Cloud backup copy (implementation) AND an attacker gains access
   & Attacker reads the cycle history
   & Exposure damages reputation (\texttt{H-REP}) or causes distress
     (\texttt{H-PSY}) \\
\bottomrule
\end{tabularx}
\end{table}

\noindent\textbf{In plain words.} Read each row from left to right: if all
conditions hold, someone gains a capability, and that capability can lead to
the harm. A protection measure blocks a pathway only if it removes one of the
pathway's \emph{necessary} conditions. Three examples:
\begin{itemize}
  \item An app lock removes the access condition of F1, so it can block F1
    even though the cycle data itself is core and stays on the phone.
  \item Turning off backup removes one branch of F2 (the server copy), but
    not the other (a seized phone). F2 stays open.
  \item No measure can remove the logged cycles without stopping the app from
    working. Pathways can only be blocked through their other conditions.
\end{itemize}
The origins are argued from the literature, not proven. The tree describes
\emph{possible} harms; it says nothing about what CycleMate actually does.

\subsection*{Step 1: Which pathways apply to CycleMate?}

Before coding, the author fills in an \emph{applicability record} for each
pathway (\Cref{tab:we-applicability}). It notes the feature the pathway needs,
where that feature is documented, the assumptions made, and a decision:
applicable, not applicable, or uncertain.

\begin{table}[H]
\centering
\caption{Applicability records for CycleMate (U1--U10 are the policy units
in Step~2).}
\label{tab:we-applicability}
\small
\begin{tabularx}{\textwidth}{@{}l>{\raggedright\arraybackslash}p{3.2cm}
                                   >{\raggedright\arraybackslash}p{2.0cm}
                                   >{\raggedright\arraybackslash}X
                                   >{\raggedright\arraybackslash}p{1.7cm}@{}}
\toprule
\textbf{ID} & \textbf{Required feature} & \textbf{Evidence} &
  \textbf{Assumptions} & \textbf{Decision} \\
\midrule
F1 & Cycle records stored on the phone & U1, U3
   & Another person can sometimes use the unlocked phone & applicable \\
F2 & Stored records, plus backup or the phone itself & U1, U3
   & A legal setting where reproductive decisions are prosecuted. This is not
     the German baseline, so it remains an assumption & uncertain \\
F3 & Identifiers shared with an analytics partner & U2
   & The partner can link the identifier across apps & applicable \\
F4 & Predictions conditional on data sharing & store listing (\enquote{premium
     predictions})
   & Not confirmed: the walkthrough has not yet been done & uncertain \\
F5 & A server copy of the cycle history & U3 (optional backup)
   & An attacker could reach the server & applicable \\
\bottomrule
\end{tabularx}
\end{table}

\noindent\textbf{In plain words.} A harm category applies to the app if at
least one of its pathways does. Psychological harm applies (F1, F5),
discrimination applies (F3), reputational harm applies (F5), and so does the
relationship tag (F1). Autonomy harm depends only on F2 and F4, which are both
uncertain, so it is \emph{uncertain}. Uncertain cases are counted and
reported separately; they are never silently dropped.

\subsection*{Step 2: Splitting the policy into units}

The policy is split at paragraph boundaries. A list stays with the sentence
that introduces it. CycleMate's policy gives ten units:

\begin{description}[leftmargin=2.4em,style=sameline,itemsep=2pt]
  \item[U1] \enquote{We collect the information you enter, such as period
    dates, symptoms, sexual activity and mood, to predict your cycle and
    fertile window.}
  \item[U2] \enquote{We share device identifiers and usage data with analytics
    partners to improve our services. We never sell your health data.}
  \item[U3] \enquote{Your health data stays on your device unless you turn on
    backup. This keeps your reproductive health private and protects you from
    anyone using it against you.}
  \item[U4] \enquote{In the unlikely event of a security incident, we will
    notify affected users as required by law.}
  \item[U5] \enquote{If your data were disclosed, it could reveal sensitive
    information such as a pregnancy, which in some jurisdictions could expose
    you to legal risk.}
  \item[U6] \enquote{In the unlikely event that an unauthorised party accessed
    your backup, you might experience unwanted contact.}
  \item[U7] \enquote{Some users worry that a partner could read their entries
    and feel hurt. There is no reason for such concern.}
  \item[U8] \enquote{If you share your cycle calendar with someone, you alone
    are responsible for any embarrassment this causes.}
  \item[U9] \enquote{Because all predictions are calculated with encryption,
    no one can ever use your data to discriminate against you.}
  \item[U10] \enquote{Your data is safe with us. We use industry-standard
    security.}
\end{description}

\subsection*{Step 3: Coding each unit}

\Cref{tab:we-units} gives the codes. In column~A, every mentioned element is
recorded with its status: \emph{asserted} (stated as fact), \emph{possible},
or \emph{denied}. Column~B is filled only when a harm is named, and the stance
is attached to that harm. Column~C records the measure, its scope and how it
is framed.

\begin{table}[H]
\centering
\caption{Unit-level codes for the CycleMate policy (codebook v3).}
\label{tab:we-units}
\small
\begin{tabularx}{\textwidth}{@{}l>{\raggedright\arraybackslash}X
                                   >{\raggedright\arraybackslash}p{2.9cm}
                                   >{\raggedright\arraybackslash}X@{}}
\toprule
\textbf{Unit} & \textbf{A: What is mentioned} & \textbf{B: Stance} &
  \textbf{C: Protection} \\
\midrule
U1 & entered data (\texttt{L-CMP}, asserted); fertile-window prediction
     (\texttt{K-INF}, asserted, as a feature) & --- & --- \\
U2 & identifiers and usage data to analytics partners (\texttt{L-CMP},
     asserted) & --- & no-sale pledge (\texttt{N-COM}); scope: health data;
     data-flow framing (\texttt{F-FLW}) \\
U3 & data on the device, optional backup (\texttt{L-CMP}); \enquote{anyone
     using it against you} (\texttt{L-GEN}, possible) & ---
   & on-device storage (\texttt{T-DEV}); scope: \enquote{anyone}, unless
     backup is on; harm framing (\texttt{F-HRM}), absolute \\
U4 & notice \enquote{as required by law} (\texttt{P-NOT}) & --- & --- \\
U5 & pregnancy inferred from disclosed data (\texttt{K-INF}, possible); legal
     risk (\texttt{H-AUT}, possible); breach consequence (\texttt{P-BRC})
   & \texttt{H-AUT}: \texttt{S-ACK} & --- \\
U6 & backup (\texttt{L-CMP}); unwanted contact (\texttt{H-PSY}, possible);
     breach consequence (\texttt{P-BRC})
   & \texttt{H-PSY}: \texttt{S-ACK} + \texttt{S-QUA} (target: likelihood of
     the breach) & --- \\
U7 & partner reads entries, user feels hurt (\texttt{H-PSY}, possible;
     \texttt{T-REL}) & \texttt{H-PSY}: \texttt{S-MIN} & --- \\
U8 & embarrassment after sharing (\texttt{H-PSY}, possible)
   & \texttt{H-PSY}: \texttt{S-RSP} & --- \\
U9 & discrimination (\texttt{H-DSC}, denied) & \texttt{H-DSC}: \texttt{S-DEN}
   & encryption (\texttt{T-ENC}); scope: \enquote{no one}, \enquote{your
     data}; harm framing (\texttt{F-HRM}), absolute \\
U10 & --- & --- & reassurance without a specific measure (\texttt{C-REA}) \\
\bottomrule
\end{tabularx}
\end{table}

\noindent The reasoning behind each row, in plain words:

\begin{description}[leftmargin=0pt,itemsep=3pt]
  \item[U1.] The user enters dates; the app \emph{derives} the fertile window.
    Deriving something the user did not enter is inference (\texttt{K-INF}).
    But the text presents it as a useful feature, not as a risk, and names no
    harm. So there is no stance.
  \item[U2.] The first sentence states a data flow (\texttt{L-CMP}); it never
    says that an analytics partner could build a profile, so \texttt{K-TRK} is
    not coded.
    \enquote{We never sell} is a promise about business practice, not a
    technical measure, so it is coded \texttt{N-COM}, kept separate from
    technical measures.
  \item[U3.] \enquote{Anyone using it against you} is vague: no concrete
    consequence is named, so it is a \emph{generic risk} (\texttt{L-GEN}), not
    a harm, and it gets no stance. The protection claim is on-device storage,
    presented as protecting the user from \enquote{anyone}: harm framing,
    absolute strength. Whether this is a mismatch is decided in Step~4.
  \item[U4.] Promising to notify users \enquote{as required by law} only
    restates a legal duty. It is coded as a general notification commitment
    (\texttt{P-NOT}), not as a concrete incident procedure (\texttt{P-INC}).
    \enquote{Unlikely} creates no stance here, because no harm is named.
  \item[U5.] A concrete consequence is named: legal risk (\texttt{H-AUT}).
    \enquote{In some jurisdictions} narrows \emph{where} the harm applies but
    does not call it unlikely, so the stance is acknowledged (\texttt{S-ACK}).
    Because the harm follows a disclosure, the unit also describes a breach
    consequence (\texttt{P-BRC}).
  \item[U6.] Unwanted contact is named and presented as possible
    (\texttt{S-ACK}). \enquote{In the unlikely event} adds likelihood language,
    so it is also qualified (\texttt{S-QUA}). The qualification targets the
    \emph{breach} (how likely it is), not the \emph{harm} (how bad it would
    be). One harm can carry two stances.
  \item[U7.] The harm (feeling hurt when a partner reads one's entries) is
    named, so it is coded \texttt{H-PSY} with the relationship tag
    (\texttt{T-REL}). It is then dismissed: \enquote{There is no reason for
    such concern} gives no reason in this unit or next to it. That makes it
    minimised (\texttt{S-MIN}). The decision rests on the wording alone, not
    on a guess about the provider's intent.
  \item[U8.] A harm is named (embarrassment), and responsibility for it is
    placed on the user (\enquote{you alone are responsible}): \texttt{S-RSP}.
  \item[U9.] Discrimination is named only to be ruled out (\enquote{no one can
    ever}): its status is \emph{denied}, and the stance is \texttt{S-DEN}. A
    measure is claimed too, so it is also coded under~C.
  \item[U10.] \enquote{Safe with us} and \enquote{industry-standard} name no
    mechanism that could be checked: reassurance without substance
    (\texttt{C-REA}). There is no measure to assess further.
\end{description}

\subsection*{Step 4: The author checks the protection claims}

Each protection claim is compared with the pathways that fall within its
scope (\Cref{tab:we-mismatch}). A \emph{potential mismatch} needs two things:
a pathway that the claim itself covers, and a written reason why the measure
does not remove one of that pathway's necessary conditions. A risk that lies
outside what the claim promises is not enough.

\begin{table}[H]
\centering
\caption{The author's review of CycleMate's protection claims.}
\label{tab:we-mismatch}
\small
\begin{tabularx}{\textwidth}{@{}l>{\raggedright\arraybackslash}p{3.4cm}
                                   >{\raggedright\arraybackslash}X
                                   >{\raggedright\arraybackslash}p{2.3cm}@{}}
\toprule
\textbf{Unit} & \textbf{What the claim covers} & \textbf{Reasoning} &
  \textbf{Decision} \\
\midrule
U2 & Sale of health data
   & F3 runs through shared identifiers and usage data, not through selling
     health data. The residual risk lies outside the claim's scope.
   & no mismatch \\
\addlinespace
U3 & Protection from \enquote{anyone}, while backup is off
   & F1 is within scope: a person using the unlocked phone is
     \enquote{anyone}. Keeping data on the phone does not remove that access
     condition. F2's seized-phone branch is also within scope, but F2 is
     uncertain (Step~1), so it is noted, not counted. F5 is outside scope,
     because the claim excludes backup.
   & potential mismatch (\texttt{C-MIS}) for F1 \\
\addlinespace
U9 & No one can use \enquote{your data} to discriminate
   & F3 is within scope: it leads to discrimination using the identifiers and
     usage data that U2 says are shared. Encrypting the predictions does not
     remove that condition.
   & potential mismatch (\texttt{C-MIS}) for F3 \\
\addlinespace
U10 & Nothing specific & No measure is named, so there is nothing to compare.
   & not assessable \\
\bottomrule
\end{tabularx}
\end{table}

\noindent A potential mismatch is a statement about the \emph{text}: what it
promises compared with what the named measure can do. It is not a finding
that CycleMate's actual system is unsafe. The study does not test any app's
security.

\subsection*{Step 5: The result for this policy}

The unit codes are now combined per app and set against the applicability
records from Step~1 (\Cref{tab:we-app}).

\begin{table}[H]
\centering
\caption{CycleMate's privacy policy against the reference tree.}
\label{tab:we-app}
\small
\begin{tabularx}{\textwidth}{@{}>{\raggedright\arraybackslash}p{3.0cm}
                                 >{\raggedright\arraybackslash}p{2.6cm}
                                 >{\raggedright\arraybackslash}X@{}}
\toprule
\textbf{Element} & \textbf{Applicable?} & \textbf{Mentioned in the policy?} \\
\midrule
Psychological (\texttt{H-PSY}) & yes (F1, F5)
  & yes: U6 (acknowledged and qualified), U7 (minimised), U8 (responsibility
    shifted) \\
Relationship tag (\texttt{T-REL}) & yes (F1) & yes: U7 (minimised) \\
Discrimination (\texttt{H-DSC}) & yes (F3) & only as denied: U9 \\
Reputational (\texttt{H-REP}) & yes (F5)
  & no mention identified in the captured policy \\
Autonomy (\texttt{H-AUT}) & uncertain (F2, F4)
  & yes: U5 (acknowledged); reported separately \\
\addlinespace
Inference (\texttt{K-INF}) & yes (F1) & yes: U1 (as a feature), U5 (as a risk) \\
Tracking (\texttt{K-TRK}) & yes (F3) & no: U2 names the data flow only \\
Breach consequence (\texttt{P-BRC}) & yes (F5) & yes: U5, U6 \\
Incident procedure (\texttt{P-INC}) & yes (F5)
  & no concrete procedure; only the legal-notice promise (U4) \\
\bottomrule
\end{tabularx}
\end{table}

\noindent What this one policy contributes to the research questions:
\begin{itemize}
  \item \textbf{RQ1 (coverage).} Four harm categories, including the
    relationship tag, apply to CycleMate. Two are named as possible harms.
    Discrimination is named only in order to be denied. Reputational harm has
    no mention in the captured policy. Autonomy harm is named, but because its
    applicability is uncertain, it is reported separately and neither counted
    as covered nor as missing.
  \item \textbf{RQ2 (stance).} Five named harms carry stances: one
    acknowledged, one acknowledged and qualified, one minimised, one with
    responsibility shifted to the user, and one denied.
  \item \textbf{RQ3 (protection).} Three protection claims are made: a
    no-sale pledge (a practice commitment), on-device storage and encryption.
    The two technical measures are framed as harm prevention with absolute
    strength, and the author flags both as potential mismatches. The no-sale
    pledge is narrowly scoped, and no reference pathway falls within its
    scope. One unit offers reassurance without substance. None of this
    verifies whether the claims are true.
  \item \textbf{RQ4 (surfaces).} The same steps are repeated for the store
    listing, the website and, for the twelve walkthrough apps, the in-app
    screens. The code sets are then compared. For example: does the legal-risk
    warning in U5 also appear on the onboarding screen where the user decides
    whether to sign up?
\end{itemize}
\enquote{No mention identified} always refers to the captured material and to
an instrument with measured limits; it is never a claim that the provider
says nothing anywhere.

\subsection*{Step 6: What the language model does}

The language model helps with Step~3 only, in three stages: it first finds
the relevant spans, then assigns a stance to each named harm, then describes
any protection claim. For every code it must quote the supporting span and
give a short reason; a code without a span is discarded. For U6, the stance
stage returns:

\begin{lstlisting}[style=jsoncode]
{"unit": "U6",
 "harm": {"code": "H-PSY",
          "span": "you might experience unwanted contact"},
 "stances": [
  {"code": "S-ACK",
   "span": "you might experience unwanted contact",
   "reason": "the harm is presented as possible"},
  {"code": "S-QUA",
   "span": "In the unlikely event that an unauthorised party accessed your backup",
   "target": "likelihood of the breach",
   "reason": "likelihood language about the event, not about the harm"}]}
\end{lstlisting}

The model never decides applicability (Step~1) or potential mismatches
(Step~4); those are the author's decisions, with the reasoning logged. Every
harm mention the model returns is checked by hand, so the model works as a
screener and is tuned to miss as little as possible. To measure how good it
is, the author codes ten validation apps completely by hand before looking at
the model's output. A detection counts as correct only if the label, the unit
and the evidence span all match. If the first stage misses a span, the codes
that depend on it count as missed as well.



\chapter{LLM Prompts}
\label{app:prompts}
\todo{Every prompt version used, verbatim, with its version number and date.}



\chapter{App Sample}
\label{app:app-sample}
This appendix lists the sample and the evidence for the developer rule
(Section~\ref{sec:app-sampling}). Tables~\ref{tab:app-sample}
to~\ref{tab:developer-rule} follow on landscape pages, and
Table~\ref{tab:developer-aliases} closes the appendix.

\section{Sample}
Table~\ref{tab:app-sample} lists the 50 apps of the sample
(Section~\ref{sec:results-sample}). Numbers 1--47 are the ranked store apps in
descending order of the German-market popularity score; 48--50 are the
preinstalled apps. \emph{Listed in} gives the stores in which the app was found.
$r_i$ is the rating count of the German App Store (23 October 2026, 19:46~UTC);
$v_i$ is the count of written reviews on German Google Play (search of 24
October 2026, 15:07~UTC; product pages of 24 and 25 October); $P_i=v_i+c\,r_i$
with $c=0.0896$ is the popularity score, rounded. A dash means that no listing
from that store is in the search frame.
Categories follow the domain of Section~\ref{sec:app-sampling}. The rows are
generated by \texttt{scripts/app-ranking/v004/rank.py} from the collections in
\texttt{05\_Data/raw/app-sample\_2026-10-25/}.

\section{Developer Rule}
\label{app:developer-rule}
In Tables~\ref{tab:developer-pilot} and~\ref{tab:developer-rule}, shaded rows
belong to a developer that holds more than one of the listed places.

Table~\ref{tab:developer-pilot} lists the 47 Google Play apps with the most
ratings in the trial collection of 22 October 2026 (search and product pages
read between 20:58 and 21:12~UTC; 19 keywords, German storefront, first 20
results; Samsung Health and apps with fewer than 1{,}000 ratings removed). That
collection ranked apps by the worldwide Play rating count and searched Google
Play only, so it has no App Store count $r_i$ and no score $P_i$; $v_i$ is the
German review count of the same pages. The privacy policy is the link on the
listing without its query string. All ten Leap Fitness Group listings link to
\path{leap.app/privacypolicy.html}; their links differ only in a
\texttt{pkg} parameter that names the app. The rows are generated by
\texttt{scripts/sampling-checks/v002/pilot\_table.py}.

Table~\ref{tab:developer-rule} repeats the selection of the sample on the same
collections and screening decisions but without the developer rule: the first
47 eligible candidates in descending order of $P_i$. Its columns are those of
Table~\ref{tab:app-sample}; \emph{Sample no.} gives the app's number in that
table, and a dash marks an app that the developer rule skips because an app of
the same provider ranks higher. The 47 places go to 37 providers. The rows are
generated by \texttt{scripts/app-ranking/v004/rank.py}.

Table~\ref{tab:developer-aliases} lists the developer names that condition P3
matched by hand, with the evidence on the listings. Each row maps the developer
name on the left to the provider on the right.

\begin{landscape}
\begin{singlespace}
\small
\setlength{\tabcolsep}{4pt}
\begin{longtable}{@{}r>{\raggedright\arraybackslash}p{5.2cm}>{\raggedright\arraybackslash}p{4.8cm}llrrr@{}}
\caption{The app sample.}\label{tab:app-sample}\\
\toprule
 & & & & & \multicolumn{2}{c}{\textbf{German count}} & \\
\cmidrule(lr){6-7}
\textbf{No.} & \textbf{App} & \textbf{Developer} & \textbf{Listed in} & \textbf{Category} & \textbf{Ratings $r_i$} & \textbf{Reviews $v_i$} & \textbf{Score $P_i$} \\
\midrule
\endfirsthead
\toprule
 & & & & & \multicolumn{2}{c}{\textbf{German count}} & \\
\cmidrule(lr){6-7}
\textbf{No.} & \textbf{App} & \textbf{Developer} & \textbf{Listed in} & \textbf{Category} & \textbf{Ratings $r_i$} & \textbf{Reviews $v_i$} & \textbf{Score $P_i$} \\
\midrule
\endhead
\bottomrule
\endfoot

% Generated by scripts/app-ranking/v004/rank.py from the collections of 2026-10-23 and 2026-10-24; do not edit by hand.
1 & komoot - Wandern und Radfahren & komoot GmbH & both & activity & 417{,}443 & 99{,}884 & 137{,}279 \\
2 & Yazio: KI Kalorienzähler & Yazio & both & diet-weight & 447{,}096 & 41{,}489 & 81{,}540 \\
3 & adidas Running: Lauftracker & adidas & both & activity & 313{,}702 & 38{,}883 & 66{,}985 \\
4 & Weight Watchers Programm & WW INTERNATIONAL, INC & both & diet-weight & 103{,}871 & 46{,}679 & 55{,}984 \\
5 & Google Health (Fitbit) & Google LLC & both & multiple & 48{,}117 & 40{,}627 & 44{,}937 \\
6 & MyFitnessPal: Kalorien Tracker & MyFitnessPal, Inc. & both & diet-weight & 80{,}204 & 32{,}562 & 39{,}747 \\
7 & Flo Zyklus- \& Periodenkalender & Flo Health Inc. & both & cycle & 218{,}522 & 14{,}357 & 33{,}932 \\
8 & Periodenkalender, Zyklus & Simple Design Ltd. & Google Play & cycle & -- & 32{,}901 & 32{,}901 \\
9 & Lifesum: Kalorienzähler AI & Lifesum & both & diet-weight & 81{,}834 & 21{,}767 & 29{,}098 \\
10 & Schrittzähler \& Pedometer & ITO Technologies, Inc. & both & activity & 46{,}124 & 18{,}692 & 22{,}824 \\
11 & Schwangerschaft + Tracker-App & Philips Electronics UK Limited & both & pregnancy & 45{,}420 & 18{,}564 & 22{,}633 \\
12 & Workouts zuhause - ohne Geräte & Leap Fitness Group & both & activity & 17{,}457 & 20{,}907 & 22{,}471 \\
13 & Clue Perioden Kalender & Clue Period Tracker by BioWink & both & cycle & 99{,}558 & 10{,}118 & 19{,}036 \\
14 & Fddb: KI Kalorienzähler \& Diät & Fddb & both & diet-weight & 17{,}446 & 15{,}600 & 17{,}163 \\
15 & Pacer Pedometer: Schrittzähler & Pacer Health & both & activity & 43{,}302 & 13{,}269 & 17{,}148 \\
16 & Freeletics: Fitness Workouts & Freeletics & both & activity & 53{,}060 & 11{,}206 & 15{,}959 \\
17 & Pregnancy App \& Baby Tracker & BabyCenter & both & pregnancy & 648 & 15{,}132 & 15{,}190 \\
18 & Kalorienzähler von fatsecret & fatsecret & both & diet-weight & 8{,}648 & 13{,}890 & 14{,}665 \\
19 & Sweatcoin Schrittzähler & Sweatco Ltd & both & activity & 30{,}005 & 11{,}695 & 14{,}383 \\
20 & StepsApp Schrittzähler & StepsApp & both & activity & 126{,}084 & 1{,}486 & 12{,}781 \\
21 & Fastic AI Abnehmen \& Fasten & HealthVida GmbH \& Co. KG. & App Store & diet-weight & 131{,}138 & -- & 11{,}747 \\
22 & Headspace Schlaf \& Meditation & Headspace for Meditation, Mindfulness and Sleep & both & mood & 84{,}651 & 4{,}141 & 11{,}724 \\
23 & Strava: Laufen \& Radfahren & Strava Inc. & both & activity & 27{,}351 & 8{,}425 & 10{,}875 \\
24 & Überwache dein Gewicht & Husain Al-Bustan & both & diet-weight & 12{,}260 & 9{,}164 & 10{,}262 \\
25 & Nike Run Club: Laufen \& Cardio & Nike, Inc. & both & activity & 40{,}450 & 6{,}305 & 9{,}929 \\
26 & Alarmy -- Wecker und Schlaf & Sleep Tracker \& Alarm Clock by Delightroom & both & sleep & 15{,}070 & 7{,}195 & 8{,}545 \\
27 & Noom: Gesund Abnehmen & Noom, Inc. & both & diet-weight & 6{,}345 & 7{,}911 & 8{,}479 \\
28 & Daylio Tagebuch \& Gewohnheiten & Habitics & both & mood & 7{,}948 & 7{,}435 & 8{,}147 \\
29 & GPS Laufen Joggen \& Radfahren & sportractive.com & Google Play & activity & -- & 8{,}091 & 8{,}091 \\
30 & Seven: 7 Minuten Workout & Perigee AB & both & activity & 66{,}701 & 2{,}036 & 8{,}011 \\
31 & Periodenkalender - Zyklus & SimpleInnovation & both & cycle & 23{,}575 & 5{,}253 & 7{,}365 \\
32 & ASICS Runkeeper: Running App & ASICS Runner App Inc. & both & activity & 24{,}090 & 5{,}105 & 7{,}263 \\
33 & STADTRADELN & Klima-Bündnis Climate Alliance Alianza del Clima & App Store & activity & 79{,}411 & -- & 7{,}114 \\
34 & Sleep as Android: Schlafzyklen & Petr Nálevka (Urbandroid) & Google Play & sleep & -- & 7{,}044 & 7{,}044 \\
35 & Macadam, Der Schrittzähler. & MacadamApp & both & activity & 28{,}105 & 3{,}136 & 5{,}654 \\
36 & Kalorien Zähler Abnehmen Arise & A.R.I.S.E. Apps resulting in self enhancement Ltd. & App Store & diet-weight & 57{,}600 & -- & 5{,}160 \\
37 & Herzfrequenzmesser - Puls & QR Code Scanner. & Google Play & heart & -- & 4{,}740 & 4{,}740 \\
38 & Sports Tracker & Suunto & Google Play & activity & -- & 4{,}521 & 4{,}521 \\
39 & Schrittzähler - Pedometer & Hitchhike Tech & Google Play & activity & -- & 4{,}352 & 4{,}352 \\
40 & Foodvisor - KI Kalorien Zähler & Foodvisor & both & diet-weight & 6{,}287 & 3{,}752 & 4{,}315 \\
41 & Relive: 3D-Videos erstellen & Relive B.V. & both & activity & 16{,}760 & 2{,}503 & 4{,}004 \\
42 & BetterMe Well-Being Coach & BetterMe Limited & App Store & activity & 44{,}421 & -- & 3{,}979 \\
43 & Bikemap: Fahrrad Navi \& GPS & Bikemap GmbH & both & activity & 19{,}737 & 2{,}182 & 3{,}950 \\
44 & Naviki -- das Fahrrad-Navi & beemo & Google Play & activity & -- & 3{,}941 & 3{,}941 \\
45 & BodyFast: Intervallfasten KI & BodyFast GmbH & App Store & diet-weight & 43{,}488 & -- & 3{,}896 \\
46 & Schrittzähler++ & Cross Forward Consulting, LLC & App Store & activity & 43{,}328 & -- & 3{,}881 \\
47 & Sleep Monitor - Schlaftracker & Fasting APP Group & Google Play & sleep & -- & 3{,}819 & 3{,}819 \\
48 & Apple Health & Apple & preinstalled & multiple & -- & -- & -- \\
49 & Apple Fitness & Apple & preinstalled & activity & -- & -- & -- \\
50 & Samsung Health & Samsung Electronics Co.,  Ltd. & preinstalled & multiple & -- & -- & -- \\


\end{longtable}

\begin{longtable}{@{}r>{\raggedright\arraybackslash}p{4.4cm}>{\raggedright\arraybackslash}p{3.8cm}>{\raggedright\arraybackslash}p{5.8cm}rr@{}}
\caption{Google Play trial collection of 22 October 2026: the 47 apps with
the most ratings worldwide.}\label{tab:developer-pilot}\\
\toprule
\textbf{No.} & \textbf{App} & \textbf{Developer} & \textbf{Privacy policy} & \textbf{Ratings (worldwide)} & \textbf{Reviews $v_i$} \\
\midrule
\endfirsthead
\toprule
\textbf{No.} & \textbf{App} & \textbf{Developer} & \textbf{Privacy policy} & \textbf{Ratings (worldwide)} & \textbf{Reviews $v_i$} \\
\midrule
\endhead
\bottomrule
\endfoot

% Generated by scripts/sampling-checks/v002/pilot_table.py from the pilot of 2026-10-22; do not edit by hand.
1 & Flo Zyklus- \& Periodenkalender & Flo Health Inc. & \path{flo.health/privacy-policy} & 5{,}154{,}162 & 14{,}356 \\
\rowcolor{devshade}2 & Periodenkalender, Zyklus & Simple Design Ltd. & \path{pc.simpledesign.ltd/privacypolicy.html} & 4{,}896{,}400 & 32{,}900 \\
\rowcolor{devshade}3 & Workouts zuhause - ohne Geräte & Leap Fitness Group & \path{leap.app/privacypolicy.html} & 4{,}045{,}384 & 20{,}907 \\
4 & Schwangerschaft + Tracker-App & Philips Electronics UK Limited & \path{info.philips-digital.com/PrivacyNotice} & 3{,}696{,}333 & 18{,}563 \\
5 & Sweatcoin Schrittzähler & Sweatco Ltd & \path{sweatco.in/privacy} & 3{,}151{,}098 & 11{,}695 \\
6 & MyFitnessPal: Kalorien Tracker & MyFitnessPal, Inc. & \path{myfitnesspal.com/privacy-policy} & 2{,}919{,}707 & 32{,}560 \\
\rowcolor{devshade}7 & Sixpack in 30 Tagen & Leap Fitness Group & \path{leap.app/privacypolicy.html} & 2{,}858{,}479 & 19{,}202 \\
\rowcolor{devshade}8 & Abnehmen für Männer-30 Tagen & Leap Fitness Group & \path{leap.app/privacypolicy.html} & 2{,}769{,}837 & 16{,}124 \\
9 & adidas Running: Lauftracker & adidas & \path{runtastic.com/privacy-notice} & 1{,}742{,}328 & 38{,}883 \\
\rowcolor{devshade}10 & Schrittzähler - Pedometer & Leap Fitness Group & \path{leap.app/privacypolicy.html} & 1{,}718{,}855 & 16{,}553 \\
11 & Mi Fitness (Xiaomi Wear) & Beijing Xiaomi Mobile Software Co.,Ltd & \path{us.watch.iot.mi.com/html/privacy/index.html} & 1{,}645{,}752 & 19{,}659 \\
12 & Pregnancy App \& Baby Tracker & BabyCenter & \path{babycenter.com/help-privacy} & 1{,}540{,}215 & 15{,}132 \\
13 & Clue Perioden Kalender & Clue Period Tracker by BioWink & \path{helloclue.com/privacy.html} & 1{,}503{,}611 & 10{,}118 \\
14 & Zepp & Zepp, Inc. & \path{upload-cdn.huami.com/tposts/8191} & 1{,}420{,}968 & 28{,}913 \\
\rowcolor{devshade}15 & Google Health (Fitbit) & Google LLC & \path{fitbit.com/privacy} & 1{,}322{,}349 & 40{,}621 \\
\rowcolor{devshade}16 & Abnehmen in 30 Tagen - Fitness & Simple Design Ltd. & \path{simpledesign.ltd/privacypolicy.html} & 1{,}249{,}774 & 6{,}907 \\
\rowcolor{devshade}17 & Abnehm-App für Frauen & Leap Fitness Group & \path{leap.app/privacypolicy.html} & 1{,}230{,}962 & 5{,}570 \\
18 & Strava: Laufen \& Radfahren & Strava Inc. & \path{strava.com/privacy} & 1{,}204{,}262 & 8{,}423 \\
19 & Garmin Connect & Garmin & \path{garmin.com/privacy/connect} & 1{,}113{,}833 & 35{,}790 \\
\rowcolor{devshade}20 & Nike Run Club: Laufen \& Cardio & Nike, Inc. & \path{nike.com/help/a/privacy-global} & 1{,}113{,}656 & 6{,}305 \\
21 & Schrittzähler \& Pedometer & ITO Technologies, Inc. & \path{ito-technologies.com/privacy} & 1{,}048{,}380 & 18{,}690 \\
22 & Pacer Pedometer: Schrittzähler & Pacer Health & \path{mypacer.com/privacy} & 977{,}049 & 13{,}270 \\
\rowcolor{devshade}23 & Schrittzähler - Pedometer & Simple Design Ltd. & \path{simpledesign.ltd/privacypolicy.html} & 956{,}554 & 9{,}421 \\
24 & Yazio: KI Kalorienzähler & Yazio & \path{help.yazio.com/hc/articles/203444951} & 873{,}216 & 41{,}477 \\
25 & Schrittzähler - Pedometer & Hitchhike Tech & \path{stepcounter.brilliantapp.ltd/privacypolicy.html} & 829{,}648 & 4{,}351 \\
\rowcolor{devshade}26 & Schrittzähler - Health , iStep & Leap Fitness Group & \path{leap.app/privacypolicy.html} & 816{,}548 & 12{,}354 \\
\rowcolor{devshade}27 & Google Fit: Aktivitätstracker & Google LLC & \path{google.com/policies/privacy} & 686{,}717 & 11{,}764 \\
28 & ASICS Runkeeper: Running App & ASICS Runner App Inc. & \path{runkeeper.com/privacypolicy} & 664{,}373 & 5{,}105 \\
29 & Calm: Meditation und Schlaf & Calm.com, Inc. & \path{calm.com/privacy} & 616{,}593 & 4{,}560 \\
30 & Periodenkalender - Zyklus & SimpleInnovation & \path{simpleinnovation.us/my-calendar/privacy-policy} & 559{,}318 & 5{,}253 \\
\rowcolor{devshade}31 & 30 Tage Fitness Challenge & Leap Fitness Group & \path{leap.app/privacypolicy.html} & 550{,}094 & 3{,}852 \\
32 & Kalorienzähler von fatsecret & fatsecret & \path{fatsecret.com/privacy-policy} & 548{,}035 & 13{,}888 \\
\rowcolor{devshade}33 & Bauchfett Verlieren & Leap Fitness Group & \path{leap.app/privacypolicy.html} & 518{,}203 & 4{,}314 \\
34 & Schwangerschafts-App & Amila & \path{amila.io/apps/privacy/privacy.html} & 476{,}342 & 565 \\
35 & Abnehmen für Frauen: Zuhause & EZ Health & \path{ezhealth.app/privacypolicy.html} & 473{,}277 & 3{,}113 \\
\rowcolor{devshade}36 & Menstruationskalender, Zyklus & Leap Fitness Group & \path{leap.app/privacypolicy.html} & 455{,}824 & 2{,}718 \\
37 & BitePal: KI-Kalorienzähler & RefaceLithuania & \path{bitepal.app/privacy} & 444{,}630 & 600 \\
38 & Da Fit & MO YOUNG LIMITED & \path{cdn.moyoung.com/HTML/dafit_law_and_privacy_en_0130.html} & 410{,}185 & 5{,}732 \\
39 & BetterSleep: Einschlafhilfe & BetterSleep Health & \path{bettersleep.com/legal/privacy-policy} & 397{,}213 & 3{,}193 \\
40 & Sleep as Android: Schlafzyklen & Petr Nálevka (Urbandroid) & \path{urbandroid.static.domains/policy.html} & 392{,}762 & 7{,}042 \\
41 & Herzfrequenzmesser - Puls & QR Code Scanner. & \path{heartrate.deepthought.industries/privacypolicy.html} & 381{,}102 & 4{,}738 \\
42 & Lifesum: Kalorienzähler AI & Lifesum & \path{lifesum.com/policy} & 378{,}999 & 21{,}767 \\
\rowcolor{devshade}43 & Nike Training Club & Nike, Inc. & \path{nike.com/help/a/privacy-global} & 372{,}660 & 2{,}587 \\
44 & Map My Walk: Walking Tracker & MapMyFitness, Inc. & \path{outsideinc.com/privacy-policy} & 368{,}932 & 725 \\
45 & Period Tracker & GP International LLC & \path{gpapps.com/support/privacy-policy} & 357{,}842 & 1{,}427 \\
46 & Clover - Periodenkalender & Wachanga & \path{wachanga.com/eng/privacy} & 314{,}039 & 1{,}338 \\
\rowcolor{devshade}47 & Fitness App: Workout zuhause & Leap Fitness Group & \path{leap.app/privacypolicy.html} & 309{,}727 & 1{,}822 \\


\end{longtable}

\setlength{\tabcolsep}{4pt}
\begin{longtable}{@{}r>{\raggedright\arraybackslash}p{4.4cm}>{\raggedright\arraybackslash}p{3.8cm}llrrrr@{}}
\caption{The selection without the developer rule (collections of 23 and 24
October 2026).}\label{tab:developer-rule}\\
\toprule
 & & & & & \multicolumn{2}{c}{\textbf{German count}} & & \\
\cmidrule(lr){6-7}
\textbf{No.} & \textbf{App} & \textbf{Developer} & \textbf{Listed in} & \textbf{Category} & \textbf{Ratings $r_i$} & \textbf{Reviews $v_i$} & \textbf{Score $P_i$} & \textbf{Sample no.} \\
\midrule
\endfirsthead
\toprule
 & & & & & \multicolumn{2}{c}{\textbf{German count}} & & \\
\cmidrule(lr){6-7}
\textbf{No.} & \textbf{App} & \textbf{Developer} & \textbf{Listed in} & \textbf{Category} & \textbf{Ratings $r_i$} & \textbf{Reviews $v_i$} & \textbf{Score $P_i$} & \textbf{Sample no.} \\
\midrule
\endhead
\bottomrule
\endfoot

% Generated by scripts/app-ranking/v004/rank.py from the collections of 2026-10-23 and 2026-10-24; do not edit by hand.
1 & komoot - Wandern und Radfahren & komoot GmbH & both & activity & 417{,}443 & 99{,}884 & 137{,}279 & 1 \\
2 & Yazio: KI Kalorienzähler & Yazio & both & diet-weight & 447{,}096 & 41{,}489 & 81{,}540 & 2 \\
3 & adidas Running: Lauftracker & adidas & both & activity & 313{,}702 & 38{,}883 & 66{,}985 & 3 \\
4 & Weight Watchers Programm & WW INTERNATIONAL, INC & both & diet-weight & 103{,}871 & 46{,}679 & 55{,}984 & 4 \\
\rowcolor{devshade}5 & Google Health (Fitbit) & Google LLC & both & multiple & 48{,}117 & 40{,}627 & 44{,}937 & 5 \\
6 & MyFitnessPal: Kalorien Tracker & MyFitnessPal, Inc. & both & diet-weight & 80{,}204 & 32{,}562 & 39{,}747 & 6 \\
7 & Flo Zyklus- \& Periodenkalender & Flo Health Inc. & both & cycle & 218{,}522 & 14{,}357 & 33{,}932 & 7 \\
\rowcolor{devshade}8 & Periodenkalender, Zyklus & Simple Design Ltd. & Google Play & cycle & -- & 32{,}901 & 32{,}901 & 8 \\
9 & Lifesum: Kalorienzähler AI & Lifesum & both & diet-weight & 81{,}834 & 21{,}767 & 29{,}098 & 9 \\
10 & Schrittzähler \& Pedometer & ITO Technologies, Inc. & both & activity & 46{,}124 & 18{,}692 & 22{,}824 & 10 \\
11 & Schwangerschaft + Tracker-App & Philips Electronics UK Limited & both & pregnancy & 45{,}420 & 18{,}564 & 22{,}633 & 11 \\
\rowcolor{devshade}12 & Workouts zuhause - ohne Geräte & Leap Fitness Group & both & activity & 17{,}457 & 20{,}907 & 22{,}471 & 12 \\
\rowcolor{devshade}13 & Sixpack in 30 Tagen & Leap Fitness Group & both & activity & 7{,}206 & 19{,}202 & 19{,}848 & -- \\
14 & Clue Perioden Kalender & Clue Period Tracker by BioWink & both & cycle & 99{,}558 & 10{,}118 & 19{,}036 & 13 \\
15 & Fddb: KI Kalorienzähler \& Diät & Fddb & both & diet-weight & 17{,}446 & 15{,}600 & 17{,}163 & 14 \\
16 & Pacer Pedometer: Schrittzähler & Pacer Health & both & activity & 43{,}302 & 13{,}269 & 17{,}148 & 15 \\
\rowcolor{devshade}17 & Schrittzähler - Pedometer & Leap Fitness Group & Google Play & activity & -- & 16{,}558 & 16{,}558 & -- \\
\rowcolor{devshade}18 & Abnehmen für Männer-30 Tagen & Leap Fitness Group & both & activity & 619 & 16{,}125 & 16{,}180 & -- \\
19 & Freeletics: Fitness Workouts & Freeletics & both & activity & 53{,}060 & 11{,}206 & 15{,}959 & 16 \\
20 & Pregnancy App \& Baby Tracker & BabyCenter & both & pregnancy & 648 & 15{,}132 & 15{,}190 & 17 \\
21 & Kalorienzähler von fatsecret & fatsecret & both & diet-weight & 8{,}648 & 13{,}890 & 14{,}665 & 18 \\
22 & Sweatcoin Schrittzähler & Sweatco Ltd & both & activity & 30{,}005 & 11{,}695 & 14{,}383 & 19 \\
23 & StepsApp Schrittzähler & StepsApp & both & activity & 126{,}084 & 1{,}486 & 12{,}781 & 20 \\
\rowcolor{devshade}24 & Schrittzähler - Health , iStep & Leap Fitness Group & Google Play & activity & -- & 12{,}354 & 12{,}354 & -- \\
\rowcolor{devshade}25 & Google Fit: Aktivitätstracker & Google LLC & Google Play & activity & -- & 11{,}765 & 11{,}765 & -- \\
26 & Fastic AI Abnehmen \& Fasten & HealthVida GmbH \& Co. KG. & App Store & diet-weight & 131{,}138 & -- & 11{,}747 & 21 \\
27 & Headspace Schlaf \& Meditation & Headspace for Meditation, Mindfulness and Sleep & both & mood & 84{,}651 & 4{,}141 & 11{,}724 & 22 \\
28 & Strava: Laufen \& Radfahren & Strava Inc. & both & activity & 27{,}351 & 8{,}425 & 10{,}875 & 23 \\
29 & Überwache dein Gewicht & Husain Al-Bustan & both & diet-weight & 12{,}260 & 9{,}164 & 10{,}262 & 24 \\
\rowcolor{devshade}30 & Nike Run Club: Laufen \& Cardio & Nike, Inc. & both & activity & 40{,}450 & 6{,}305 & 9{,}929 & 25 \\
\rowcolor{devshade}31 & Schrittzähler - Pedometer & Simple Design Ltd. & Google Play & activity & -- & 9{,}421 & 9{,}421 & -- \\
32 & Alarmy -- Wecker und Schlaf & Sleep Tracker \& Alarm Clock by Delightroom & both & sleep & 15{,}070 & 7{,}195 & 8{,}545 & 26 \\
33 & Noom: Gesund Abnehmen & Noom, Inc. & both & diet-weight & 6{,}345 & 7{,}911 & 8{,}479 & 27 \\
34 & Daylio Tagebuch \& Gewohnheiten & Habitics & both & mood & 7{,}948 & 7{,}435 & 8{,}147 & 28 \\
35 & GPS Laufen Joggen \& Radfahren & sportractive.com & Google Play & activity & -- & 8{,}091 & 8{,}091 & 29 \\
36 & Seven: 7 Minuten Workout & Perigee AB & both & activity & 66{,}701 & 2{,}036 & 8{,}011 & 30 \\
37 & Periodenkalender - Zyklus & SimpleInnovation & both & cycle & 23{,}575 & 5{,}253 & 7{,}365 & 31 \\
38 & ASICS Runkeeper: Running App & ASICS Runner App Inc. & both & activity & 24{,}090 & 5{,}105 & 7{,}263 & 32 \\
39 & STADTRADELN & Klima-Bündnis Climate Alliance Alianza del Clima & App Store & activity & 79{,}411 & -- & 7{,}114 & 33 \\
40 & Sleep as Android: Schlafzyklen & Petr Nálevka (Urbandroid) & Google Play & sleep & -- & 7{,}044 & 7{,}044 & 34 \\
\rowcolor{devshade}41 & Abnehmen in 30 Tagen - Fitness & Simple Design Ltd. & Google Play & activity & -- & 6{,}908 & 6{,}908 & -- \\
\rowcolor{devshade}42 & Nike Training Club & Nike, Inc. & both & activity & 36{,}896 & 2{,}587 & 5{,}892 & -- \\
43 & Macadam, Der Schrittzähler. & MacadamApp & both & activity & 28{,}105 & 3{,}136 & 5{,}654 & 35 \\
\rowcolor{devshade}44 & Abnehm-App für Frauen & Leap Fitness Group & Google Play & activity & -- & 5{,}569 & 5{,}569 & -- \\
\rowcolor{devshade}45 & Intervallfasten - Fasten-App & Leap Fitness Group & both & diet-weight & 2{,}650 & 5{,}175 & 5{,}412 & -- \\
46 & Kalorien Zähler Abnehmen Arise & A.R.I.S.E. Apps resulting in self enhancement Ltd. & App Store & diet-weight & 57{,}600 & -- & 5{,}160 & 36 \\
47 & Herzfrequenzmesser - Puls & QR Code Scanner. & Google Play & heart & -- & 4{,}740 & 4{,}740 & 37 \\


\end{longtable}
\end{singlespace}
\end{landscape}

\begin{singlespace}
\small
\setlength{\tabcolsep}{4pt}
\begin{longtable}{@{}>{\raggedright\arraybackslash}p{3.9cm}>{\raggedright\arraybackslash}p{3.9cm}>{\raggedright\arraybackslash}p{5.9cm}@{}}
\caption{Developer names matched by hand (condition P3).}\label{tab:developer-aliases}\\
\toprule
\textbf{Developer as listed} & \textbf{Counted as} & \textbf{Evidence} \\
\midrule
\endfirsthead
\toprule
\textbf{Developer as listed} & \textbf{Counted as} & \textbf{Evidence} \\
\midrule
\endhead
\bottomrule
\endfoot

% Generated by scripts/app-ranking/v004/rank.py from the collections of 2026-10-23 and 2026-10-24; do not edit by hand.
Everyday Health, Inc. (App Store) & BabyCenter (Google Play) & App Store seller URL babycenter.com is the Play developer website \\
Clue Period Tracker by BioWink (Google Play) & BioWink GmbH (App Store) & App Store developer BioWink GmbH; both stores link helloclue.com \\
Sleep Tracker \& Alarm Clock by Delightroom (Google Play) & Delight Room Co., Ltd. (App Store) & Same app (Alarmy); App Store developer Delight Room; both link alar.my \\
babelli: Schwangerschaft \& Baby (Google Play) & Fabulabs GmbH (App Store) & App Store developer Fabulabs GmbH; both link babelli.de \\
Fddb (Google Play) & Food Database GmbH (App Store) & App Store developer Food Database GmbH; both link fddb.info \\
Relaxio s.r.o. (App Store) & Habitics s.r.o. (App Store); Habitics (Google Play) & Same app (Daylio); App Store seller URL daylio.net is the Play privacy-policy domain \\
Headspace for Meditation, Mindfulness and Sleep (Google Play) & Headspace Inc. (App Store) & App Store developer Headspace Inc.; both link headspace.com \\
BetterSleep Health (Google Play) & Ipnos Software Inc. (App Store) & App Store seller BetterSleep Health, LLC; both link bettersleep.com \\
Leap Health (App Store) & Leap Fitness Group (Google Play) & App Store seller Leap Health Fitness Limited; seller URL leap.app, the domain of the Play privacy policy \\
MacadamApp (Google Play) & MACADAM TECHNOLOGIES, S.L. (App Store) & App Store developer MACADAM TECHNOLOGIES; both link macadam.app \\
Outside Television Inc (App Store) & MapMyFitness, Inc. (Google Play) & Map My Run, Map My Ride, Map My Walk; App Store seller URLs are the Play developer websites \\
Hangzhou Youzijie Information Technology Co., Ltd. (App Store) & MEETYOU PTE.LTD. (Google Play) & Same app (MeetYou); both link meetyouintl.com \\
Philips Electronics UK Limited (Google Play) & Philips Digital UK Limited (App Store) & Same app (Pregnancy+); the Play policy and website are on philips-digital.com, which the App Store account Philips Digital UK links in the description of its Baby+ listing \\


\end{longtable}
\end{singlespace}



% Migrated from proposal/proposal.tex (5 October 2026) on 14 October 2026.
\chapter{Worked Example: From a Pathway to Text Codes}
\label{app:example}

This example applies the method to a constructed privacy notice of a
fictional period tracker, \emph{App~F}; it describes no real app.

\textbf{Step 1: Model the tree.} Figure~\ref{fig:example-path} adapts
Knabenhans et al.'s targeted-advertising tree (their Fig.~6c, with the
tracking sub-tree of Fig.~5b) to App~F~\parencite[\S4.2]{knabenhans2026harms}.
The components keep the paper's roles: C1 is their \emph{reveal other
attributes}, C2 a \emph{persistent identifier}, C3 \emph{online metadata},
which the paper says includes advertising IDs, C4 \emph{repeated
observations}, and C5 the \emph{digitization of data records} that enables
data brokerage. OR means one branch suffices; OR+ means several components
combine. App~F's core function, predicting the cycle, can run on the phone and
reveals nothing to an ad partner, so both branches need implementation
choices that the provider controls and can describe. Dropping the ad SDK
blocks tracking, and not keeping or not selling server records blocks
brokerage; sending interest categories instead of entries does not, because a
category such as \enquote{family planning} still carries the
pregnancy~\parencite{martin2025cookies}.

\begin{figure}[htbp]
\centering
\resizebox{\textwidth}{!}{%
\begin{tikzpicture}[
  font=\footnotesize,
  bx/.style={draw=primary!65,fill=primary!4,rounded corners=2pt,align=center,text width=28mm,
             minimum height=17mm,inner sep=3pt},
  lb/.style={font=\scriptsize\itshape,fill=white,inner sep=1pt}]
  \node[bx,fill=primary!12,text width=100mm,minimum height=12mm] (h) at (6.6,4.5)
    {\textbf{Harm:} targeted advertising, e.g.\ pregnancy ads after a
     pregnancy loss (psychological)};
  \node[bx,fill=accent!10,text width=70mm,minimum height=12mm] (t) at (4.95,2.25) {\textbf{Capability:} tracking/profiling};
  \node[bx,fill=accent!10,text width=28mm,minimum height=12mm] (d) at (13.2,2.25) {\textbf{Capability:} data brokerage};
  \node[bx] (c1) at (0.0,0) {C1: Pregnancy Mode events\\sent to an ad SDK\\(implementation)};
  \node[bx] (c2) at (3.3,0) {C2: user ID shared\\with ad partners\\(implementation)};
  \node[bx] (c3) at (6.6,0) {C3: advertising ID,\\IP address\\(digital context)};
  \node[bx] (c4) at (9.9,0) {C4: daily use\\of the app\\(core)};
  \node[bx] (c5) at (13.2,0) {C5: cycle records on\\the provider's servers\\(implementation)};
  \foreach \c in {c1,c2,c3,c4} \draw (\c.north) -- (t.south);
  \draw (c5.north) -- (d.south);
  \draw (t.north) -- (h.south);
  \draw (d.north) -- (h.south);
  \node[lb] at (4.95,1.35) {OR+};
  \node[lb] at (6.6,3.65) {OR};
\end{tikzpicture}}
\caption{App~F's version of Knabenhans et al.'s targeted-advertising tree,
read from the bottom up.}
\Description{A three-level tree. Four components, Pregnancy Mode events sent
to an ad SDK, a user ID shared with ad partners, the advertising ID and IP
address, and daily use of the app, combine (OR+) into tracking and profiling.
Cycle records on the provider's servers enable data brokerage. Either
capability (OR) can lead to targeted advertising, such as pregnancy ads after
a pregnancy loss, a psychological harm.}
\label{fig:example-path}
\end{figure}

\textbf{Step 2: Decide whether the tree applies.} The applicability record
notes the required features (C1 and C3 reach ad partners; C5 exists), the
evidence (S2, S3 below), the assumption (the shared events include Pregnancy
Mode, which the text does not state) and the decision: \emph{applicable}, so
psychological harm applies to App~F. Mentions are coded regardless;
applicability only decides which \emph{missing} harms count as gaps (RQ1).

\textbf{Step 3: Code the statements.} Each statement in Table~\ref{tab:worked}
is one unit. A: every mentioned element with its status (asserted, possible,
denied); B: a stance, only for a named harm; C: the measure's framing,
strength and scope. A dash: not applicable.

\begin{table}[H]
\centering
\begin{singlespace}
\caption{Constructed statements and their codes.}
\label{tab:worked}
\small
\begin{tabularx}{\textwidth}{@{}L{0.5cm}>{\raggedright\arraybackslash}X>{\raggedright\arraybackslash}X L{2.4cm}>{\raggedright\arraybackslash}X@{}}
\toprule
 & \textbf{Statement} & \textbf{A: What is mentioned} & \textbf{B: Stance}
 & \textbf{C: Protection} \\
\midrule
S1 & \enquote{We use your cycle entries to predict your fertile window.}
   & Entries: component, asserted. Prediction: inference, asserted (a
     feature) & -- & -- \\
\addlinespace
S2 & \enquote{Our advertising partners receive app usage events and your
     advertising ID.}
   & C1 and C3 to partners: components, asserted & -- & -- \\
\addlinespace
S3 & \enquote{We store your entries on our servers and never sell them.
     Partners only receive interest categories, not your entries.}
   & C5: component, asserted. Sale: denied. Entries to partners: denied.
     Categories to partners: asserted
   & -- & No-sale pledge (practice commitment); inferences instead of raw
     data; data-flow framing; absolute \\
\addlinespace
S4 & \enquote{Partners never see your entries, so you will never get
     pregnancy-based ads.}
   & Entries to partners: denied. Pregnancy-based ads (profiling):
     denied & -- (no harm named)
   & Inferences instead of raw data; harm framing; absolute; scope:
     pregnancy-based ads \\
\addlinespace
S5 & \enquote{Pregnancy ads can upset users, e.g.\ after a loss, but this
     rarely happens.}
   & Distress from such ads: psychological harm, possible
   & Acknowledged and qualified (how often) & -- \\
\bottomrule
\end{tabularx}
\end{singlespace}
\end{table}

\noindent\textbf{Why these codes.} S1: a prediction made for the user gives
no one else access; no harm is named, so no stance. S2: the data flow is named
but not what partners can do with it, so tracking is not coded. S4: a
capability is denied, not a harm, so no stance. S5: \enquote{rarely}
qualifies how often the harm occurs, so it carries two stances. Dismissing a
capability without naming a harm would be reassurance, not minimisation.

\textbf{Step 4: Review the protection claims.} S3's no-sale pledge covers the
brokerage branch and, taken at its word, closes it: no mismatch. Tracking lies
outside its scope, since C1 and C3 still flow (S2); S3 leaves it open without
contradicting it, which RQ3 records as data-flow framing. S4 promises no
pregnancy-based ads, placing tracking inside its scope; withholding entries
removes neither C1 nor C3, and categories still carry the pregnancy: the
author records a \emph{potential mismatch} with this reasoning. A mismatch
concerns what the text promises, not whether App~F is safe.

\textbf{Step 5: Summarise.} Psychological harm applies and is mentioned once
(S5); tracking is never named, although S2 describes its components. RQ3: one
scoped data-flow claim without mismatch (S3), one absolute harm-framed claim
with a potential mismatch (S4).



% Source of the listings (read from the repository at compile time, not copied):
%   App Store collector   scripts/appstore-scraper/v007/    scraper.py, input.json
%   Google Play collector scripts/googleplay-scraper/v005/  scraper.py, input.json, requirements.txt
% To show a newer version, change the paths below and the two comments above.
\chapter{Scraping Scripts}
\label{app:scripts}
Both collectors read a JSON file and print the result as JSON on standard
output. They record what each store returns for a keyword search. They apply no
eligibility rule and no ranking, and they do not observe app behaviour. The App
Store collector is the version that collected the sample on 23 October 2026,
shown unchanged. The Google Play collector is a shortened form of the version
that collected the sample on 24 and 25 October 2026
(\texttt{googleplay-scraper/v004}, kept in the repository): it drops that
version's page limit, which the sample run did not reach, and gives the same
results from the same search hits and product pages. The collected files are
kept unchanged in \texttt{05\_Data/raw/app-sample\_2026-10-25/}. The inputs
hold the 49 queries of Table~\ref{tab:keywords}.

\section{Input Format}
\begin{description}
  \item[\texttt{keywords}] Search terms, one request each (both scripts).
  \item[\texttt{region}] Two-letter storefront; \texttt{de} selects the German store.
  \item[\texttt{language}] Interface language requested from Google Play.
  \item[\texttt{limit}] App Store: maximum results per keyword (up to 200).
  \item[\texttt{depth}] Google Play: search positions whose apps get their
    product page read (30).
  \item[\texttt{cache}] Google Play: file to which each product page is appended.
\end{description}

\par\medskip\noindent
\begin{minipage}{\linewidth}
\lstinputlisting[style=jsoncode, caption={Input of the App Store collector; the remaining 46 keywords of Table~\ref{tab:keywords} are left out after line~5.}, label={lst:input-apple}, linerange={1-5,51-55}]{../scripts/appstore-scraper/v007/input.json}
\end{minipage}
\par\medskip

\noindent
\begin{minipage}{\linewidth}
\lstinputlisting[style=jsoncode, caption={Input of the Google Play collector; keywords left out as above.}, label={lst:input-play}, linerange={1-5,51-57}]{../scripts/googleplay-scraper/v005/input.json}
\end{minipage}
\par\medskip

\section{App Store Collector}
The script calls Apple's documented iTunes Search API once per keyword and
storefront, keeps iPhone-compatible results in the order returned with their
position before filtering, logs each request with its URL, time, result count and
SHA-256 hash, and pauses between requests to respect the rate limit. Rating counts
and averages describe the requested storefront (Table~\ref{tab:storefront-fields}).

\lstinputlisting[style=pythoncode, caption={App Store collector.}, label={lst:scraper-apple}]{../scripts/appstore-scraper/v007/scraper.py}

\section{Google Play Collector}
The script uses the open-source \texttt{google-play-scraper} parser (version
1.2.7). It first reads the embedded data of the first search page for every
keyword and records each hit's position, title, developer and install band; the
app ID of a highlighted top result is read from its details URL, which the
parser's own field misses. It then opens the product page of every app found
within \texttt{depth}, appending each to the \texttt{cache} file so that an
interrupted run resumes. It restores TLS certificate verification, which the
library disables on import, and rejects any request or redirect that drops the
storefront parameter, so a result cannot silently come from another country.
Failed lookups are recorded under \texttt{errors}. Search positions follow the
order of the search response and are not verified as the store's native
ranking. Rating counts and install figures are worldwide on Google Play; the
one-to-five-star breakdown is the country's distribution scaled to the
worldwide count, and the count of written reviews, which the ranking uses,
follows the storefront country (Table~\ref{tab:storefront-fields}).

\lstinputlisting[style=pythoncode, caption={Google Play collector.}, label={lst:scraper-play}]{../scripts/googleplay-scraper/v005/scraper.py}



% The table rows are generated by scripts/literature-search/v001/report.py
% (tables/literature-included.tex); change the screening file or the script, not the rows.
\chapter{Literature Review Records}
\label{app:literature}
Table~\ref{tab:lit-included} lists the works included by the literature search
of Section~\ref{sec:lit-review}, sorted by inclusion criterion, year and first
author. The codes are those of the screening (I1--I4). Each work is linked by
its DOI or, where it has none, by its page in the database that found it. The
screening of all 953 records, with the code of every excluded record, is kept
in \texttt{01\_Literature/Systematic-Review/screening.csv}, and the search
results are kept unchanged in
\texttt{05\_Data/raw/literature-search\_2026-10-26/}.

{\small
\setlength{\tabcolsep}{4pt}
\begin{longtable}{@{}r l >{\raggedright\arraybackslash}p{2.1cm} r >{\raggedright\arraybackslash}p{5.3cm} >{\raggedright\arraybackslash}p{3.6cm}@{}}
\caption{Works included by the literature search.}\label{tab:lit-included}\\
\toprule
\textbf{No.} & \textbf{Code} & \textbf{First author} & \textbf{Year} & \textbf{Title} & \textbf{DOI or link} \\
\midrule
\endfirsthead
\toprule
\textbf{No.} & \textbf{Code} & \textbf{First author} & \textbf{Year} & \textbf{Title} & \textbf{DOI or link} \\
\midrule
\endhead
\bottomrule
\endfoot

% Generated by scripts/literature-search/v001/report.py from the search of 2026-10-26 and its screening; do not edit by hand.
1 & I1 & Cyr & 2014 & Security Analysis of Wearable Fitness Devices ( Fitbit ) & \href{https://www.semanticscholar.org/paper/f4abebef4e39791f358618294cd8d040d7024399}{semanticscholar.org} \\
2 & I1 & Dehling & 2015 & Exploring the Far Side of Mobile Health: Information Security and Privacy of Mobile Health Apps on iOS and Android & \href{https://doi.org/10.2196/mhealth.3672}{\nolinkurl{10.2196/mhealth.3672}} \\
3 & I1 & Liu & 2015 & When Good Becomes Evil & \href{https://doi.org/10.1145/2810103.2813668}{\nolinkurl{10.1145/2810103.2813668}} \\
4 & I1 & Nia & 2015 & Physiological Information Leakage: A New Frontier in Health Information Security & \href{https://doi.org/10.1109/tetc.2015.2478003}{\nolinkurl{10.1109/tetc.2015.2478003}} \\
5 & I1 & Zubaydi & 2015 & Security of mobile health (mHealth) systems & \href{https://doi.org/10.1109/bibe.2015.7367689}{\nolinkurl{10.1109/bibe.2015.7367689}} \\
6 & I1 & Ching & 2016 & Wearable Technology Devices Security and Privacy Vulnerability Analysis & \href{https://doi.org/10.5121/ijnsa.2016.8302}{\nolinkurl{10.5121/ijnsa.2016.8302}} \\
7 & I1 & Goyal & 2016 & Mind the tracker you wear & \href{https://doi.org/10.1145/2851613.2851685}{\nolinkurl{10.1145/2851613.2851685}} \\
8 & I1 & Lee & 2016 & Security threat on wearable services: Empirical study using a commercial smartband & \href{https://doi.org/10.1109/icce-asia.2016.7804766}{\nolinkurl{10.1109/icce-asia.2016.7804766}} \\
9 & I1 & Liu & 2016 & Smart Attacks against Intelligent Wearables in People-Centric Internet of Things & \href{https://doi.org/10.1109/mcom.2016.1600553cm}{\nolinkurl{10.1109/mcom.2016.1600553cm}} \\
10 & I1 & Aktypi & 2017 & Unwinding Ariadne's Identity Thread & \href{https://doi.org/10.1145/3137616.3137617}{\nolinkurl{10.1145/3137616.3137617}} \\
11 & I1 & Banerjee & 2017 & Wearable devices and healthcare: Data sharing and privacy & \href{https://doi.org/10.1080/01972243.2017.1391912}{\nolinkurl{10.1080/01972243.2017.1391912}} \\
12 & I1 & Crager & 2017 & Information Leakage through Mobile Motion Sensors: User Awareness and Concerns & \href{https://doi.org/10.14722/eurousec.2017.23013}{\nolinkurl{10.14722/eurousec.2017.23013}} \\
13 & I1 & Sampat & 2017 & Privacy Risks and Security Threats in mHealth apps & \href{https://doi.org/10.58729/1941-6679.1353}{\nolinkurl{10.58729/1941-6679.1353}} \\
14 & I1 & Shrestha & 2017 & An Offensive and Defensive Exposition of Wearable Computing & \href{https://doi.org/10.1145/3133837}{\nolinkurl{10.1145/3133837}} \\
15 & I1 & Classen & 2018 & Anatomy of a Vulnerable Fitness Tracking System & \href{https://doi.org/10.1145/3191737}{\nolinkurl{10.1145/3191737}} \\
16 & I1 & Datta & 2018 & A Survey of Privacy Concerns in Wearable Devices & \href{https://doi.org/10.1109/bigdata.2018.8622110}{\nolinkurl{10.1109/bigdata.2018.8622110}} \\
17 & I1 & Hassan & 2018 & Analysis of Privacy Protections in Fitness Tracking Social Networks -or- You can run, but can you hide? & \href{https://www.semanticscholar.org/paper/a8423a5040ae69c69ff8d4e32f3e406e65fe93b1}{semanticscholar.org} \\
18 & I1 & Liu & 2018 & aLeak: Privacy Leakage through Context - Free Wearable Side-Channel & \href{https://doi.org/10.1109/infocom.2018.8485958}{\nolinkurl{10.1109/infocom.2018.8485958}} \\
19 & I1 & Celosia & 2019 & Saving private addresses: an analysis of privacy issues in the bluetooth-low-energy advertising mechanism & \href{https://doi.org/10.1145/3360774.3360777}{\nolinkurl{10.1145/3360774.3360777}} \\
20 & I1 & Cilliers & 2019 & Wearable devices in healthcare: Privacy and information security issues & \href{https://doi.org/10.1177/1833358319851684}{\nolinkurl{10.1177/1833358319851684}} \\
21 & I1 & Grundy & 2019 & Data sharing practices of medicines related apps and the mobile ecosystem: traffic, content, and network analysis & \href{https://doi.org/10.1136/bmj.l920}{\nolinkurl{10.1136/bmj.l920}} \\
22 & I1 & Kröger & 2019 & Unexpected Inferences from Sensor Data: A Hidden Privacy Threat in the Internet of Things & \href{https://doi.org/10.1007/978-3-030-15651-0_13}{\nolinkurl{10.1007/978-3-030-15651-0\_13}} \\
23 & I1 & Nobakht & 2019 & PGFit: Static permission analysis of health and fitness apps in IoT programming frameworks & \href{https://doi.org/10.1016/j.jnca.2019.102509}{\nolinkurl{10.1016/j.jnca.2019.102509}} \\
24 & I1 & Schneegaß & 2019 & Understanding the Impact of Information Representation on Willingness to Share Information & \href{https://doi.org/10.1145/3290605.3300753}{\nolinkurl{10.1145/3290605.3300753}} \\
25 & I1 & Tseng & 2019 & Threat Analysis for Wearable Health Devices and Environment Monitoring Internet of Things Integration System & \href{https://doi.org/10.1109/access.2019.2946081}{\nolinkurl{10.1109/access.2019.2946081}} \\
26 & I1 & Xue & 2019 & A review on intelligent wearables: Uses and risks & \href{https://doi.org/10.1002/hbe2.173}{\nolinkurl{10.1002/hbe2.173}} \\
27 & I1 & Blow & 2020 & A Study on Vulnerabilities and Threats to Wearable Devices & \href{https://doi.org/10.53735/cisse.v7i1.113}{\nolinkurl{10.53735/cisse.v7i1.113}} \\
28 & I1 & Conti-Cook & 2020 & Surveilling the Digital Abortion Diary & \href{https://doi.org/10.2139/ssrn.3666305}{\nolinkurl{10.2139/ssrn.3666305}} \\
29 & I1 & Gabriele & 2020 & Understanding Fitness Tracker Users' Security and Privacy Knowledge, Attitudes and Behaviours & \href{https://doi.org/10.1145/3313831.3376651}{\nolinkurl{10.1145/3313831.3376651}} \\
30 & I1 & Galvin & 2020 & Developments in Privacy and Data Ownership in Mobile Health Technologies, 2016-2019 & \href{https://doi.org/10.1055/s-0040-1701987}{\nolinkurl{10.1055/s-0040-1701987}} \\
31 & I1 & Heidel & 2020 & Potential Benefits and Risks Resulting From the Introduction of Health Apps and Wearables Into the German Statutory Health Care System: Scoping Review & \href{https://doi.org/10.2196/16444}{\nolinkurl{10.2196/16444}} \\
32 & I1 & Iwaya & 2020 & Security and Privacy for mHealth and uHealth Systems: A Systematic Mapping Study & \href{https://doi.org/10.1109/access.2020.3015962}{\nolinkurl{10.1109/access.2020.3015962}} \\
33 & I1 & Marks & 2020 & Emergent Medical Data: Health Information Inferred by Artificial Intelligence & \href{https://www.semanticscholar.org/paper/6ec346394b4f3f54d4de9a4e8b35b4f96ef8f21c}{semanticscholar.org} \\
34 & I1 & Psychoula & 2020 & Privacy Risk Awareness in Wearables and the Internet of Things & \href{https://doi.org/10.1109/mprv.2020.2997616}{\nolinkurl{10.1109/mprv.2020.2997616}} \\
35 & I1 & Wang & 2020 & BlueDoor & \href{https://doi.org/10.1145/3386901.3389025}{\nolinkurl{10.1145/3386901.3389025}} \\
36 & I1 & Zhang & 2020 & Security and Privacy Analysis of Wearable Health Device & \href{https://doi.org/10.1109/compsac48688.2020.00044}{\nolinkurl{10.1109/compsac48688.2020.00044}} \\
37 & I1 & Brown & 2021 & THE FEMTECH PARADOX: HOW WORKPLACE MONITORING & \href{https://www.semanticscholar.org/paper/3656443cf5ae7e8eb17af76cbef178fda75b8978}{semanticscholar.org} \\
38 & I1 & Collins & 2021 & Is That Lawful? Data Privacy and Fitness Trackers in the Workplace & \href{https://doi.org/10.54648/ijcl2021003}{\nolinkurl{10.54648/ijcl2021003}} \\
39 & I1 & Gilman & 2021 & Periods for Profit and the Rise of Menstrual Surveillance & \href{https://doi.org/10.52214/cjgl.v41i1.8824}{\nolinkurl{10.52214/cjgl.v41i1.8824}} \\
40 & I1 & Ioannidou & 2021 & On General Data Protection Regulation Vulnerabilities and Privacy Issues, for Wearable Devices and Fitness Tracking Applications & \href{https://doi.org/10.3390/cryptography5040029}{\nolinkurl{10.3390/cryptography5040029}} \\
41 & I1 & Jiang & 2021 & Research on Data Security and Privacy Protection of Wearable Equipment in Healthcare & \href{https://doi.org/10.1155/2021/6656204}{\nolinkurl{10.1155/2021/6656204}} \\
42 & I1 & Velykoivanenko & 2021 & Are Those Steps Worth Your Privacy? & \href{https://doi.org/10.1145/3494960}{\nolinkurl{10.1145/3494960}} \\
43 & I1 & Almeida & 2022 & Bodies Like Yours: Enquiring Data Privacy in FemTech & \href{https://doi.org/10.1145/3547522.3547674}{\nolinkurl{10.1145/3547522.3547674}} \\
44 & I1 & Barua & 2022 & Security and Privacy Threats for Bluetooth Low Energy in IoT and Wearable Devices: A Comprehensive Survey & \href{https://doi.org/10.1109/ojcoms.2022.3149732}{\nolinkurl{10.1109/ojcoms.2022.3149732}} \\
45 & I1 & Dhondt & 2022 & A Run a Day Won't Keep the Hacker Away & \href{https://doi.org/10.1145/3548606.3560616}{\nolinkurl{10.1145/3548606.3560616}} \\
46 & I1 & Iwaya & 2022 & On the privacy of mental health apps & \href{https://doi.org/10.1007/s10664-022-10236-0}{\nolinkurl{10.1007/s10664-022-10236-0}} \\
47 & I1 & Kolfschooten & 2022 & The mHealth Power Paradox & \href{https://doi.org/10.1017/9781108975452.006}{\nolinkurl{10.1017/9781108975452.006}} \\
48 & I1 & Mehrnezhad & 2022 & Vision: Too Little too Late? Do the Risks of FemTech already Outweigh the Benefits? & \href{https://doi.org/10.1145/3549015.3554204}{\nolinkurl{10.1145/3549015.3554204}} \\
49 & I1 & Mink & 2022 & Users Can Deduce Sensitive Locations Protected by Privacy Zones on Fitness Tracking Apps & \href{https://doi.org/10.1145/3491102.3502136}{\nolinkurl{10.1145/3491102.3502136}} \\
50 & I1 & Sam & 2022 & The Effectiveness of IoT Based Wearable Devices and Potential Cybersecurity Risks: A Systematic Literature Review from the Last Decade & \href{https://doi.org/10.3991/ijoe.v18i09.32255}{\nolinkurl{10.3991/ijoe.v18i09.32255}} \\
51 & I1 & Fuente & 2023 & Analysis of security and privacy issues in wearables for minors & \href{https://doi.org/10.1007/s11276-022-03211-6}{\nolinkurl{10.1007/s11276-022-03211-6}} \\
52 & I1 & Locatelli & 2023 & Device discovery and tracing in the Bluetooth Low Energy domain & \href{https://doi.org/10.1016/j.comcom.2023.02.008}{\nolinkurl{10.1016/j.comcom.2023.02.008}} \\
53 & I1 & McDonald & 2023 & “I Did Watch ‘The Handmaid's Tale’”: Threat Modeling Privacy Post-roe in the United States & \href{https://doi.org/10.1145/3589960}{\nolinkurl{10.1145/3589960}} \\
54 & I1 & Mehrnezhad & 2023 & “My sex-related data is more sensitive than my financial data and I want the same level of security and privacy”: User Risk Perceptions and Protective Actions in Female-oriented Technologies & \href{https://doi.org/10.1145/3617072.3617100}{\nolinkurl{10.1145/3617072.3617100}} \\
55 & I1 & Mishra & 2023 & Femtech apps and quantification of the reproductive body in India: Issues and concerns & \href{https://doi.org/10.1177/00113921231206491}{\nolinkurl{10.1177/00113921231206491}} \\
56 & I1 & Taka & 2023 & A deep dive into dynamic data flows, wearable devices, and the concept of health data & \href{https://doi.org/10.1093/idpl/ipad007}{\nolinkurl{10.1093/idpl/ipad007}} \\
57 & I1 & Thapa & 2023 & Security Risks and User Perception towards Adopting Wearable Internet of Medical Things & \href{https://doi.org/10.3390/ijerph20085519}{\nolinkurl{10.3390/ijerph20085519}} \\
58 & I1 & Zufferey & 2023 & \enquote{Revoked just now!} Users' Behaviors Toward Fitness-Data Sharing with Third-Party Applications & \href{https://doi.org/10.56553/popets-2023-0004}{\nolinkurl{10.56553/popets-2023-0004}} \\
59 & I1 & Cao & 2024 & \enquote{I Deleted It After the Overturn of Roe v. Wade}: Understanding Women's Privacy Concerns Toward Period-Tracking Apps in the Post Roe v. Wade Era & \href{https://doi.org/10.1145/3613904.3642042}{\nolinkurl{10.1145/3613904.3642042}} \\
60 & I1 & Chen & 2024 & Private Data Leakage in Federated Human Activity Recognition for Wearable Healthcare Devices & \href{https://doi.org/10.48550/arxiv.2405.10979}{\nolinkurl{10.48550/arxiv.2405.10979}} \\
61 & I1 & Hammond & 2024 & Intimate harms and menstrual cycle tracking apps & \href{https://doi.org/10.1016/j.clsr.2024.106038}{\nolinkurl{10.1016/j.clsr.2024.106038}} \\
62 & I1 & Howe & 2024 & Ethical Risks of Systematic Menstrual Tracking in Sport & \href{https://doi.org/10.1007/s11673-023-10333-9}{\nolinkurl{10.1007/s11673-023-10333-9}} \\
63 & I1 & Kaur & 2024 & Security and Privacy of Wearable Wireless Sensors in Healthcare: A Systematic Review & \href{https://doi.org/10.37256/cnc.2120243852}{\nolinkurl{10.37256/cnc.2120243852}} \\
64 & I1 & Mehrnezhad & 2024 & Mind the FemTech gap: regulation failings and exploitative systems & \href{https://doi.org/10.3389/friot.2024.1296599}{\nolinkurl{10.3389/friot.2024.1296599}} \\
65 & I1 & Niksirat & 2024 & Wearable Activity Trackers: A Survey on Utility, Privacy, and Security & \href{https://doi.org/10.1145/3645091}{\nolinkurl{10.1145/3645091}} \\
66 & I1 & Sivakumar & 2024 & Addressing privacy concerns with wearable health monitoring technology & \href{https://doi.org/10.1002/widm.1535}{\nolinkurl{10.1002/widm.1535}} \\
67 & I1 & Yeke & 2024 & Wear’s my Data? Understanding the Cross-Device Runtime Permission Model in Wearables & \href{https://doi.org/10.1109/sp54263.2024.00077}{\nolinkurl{10.1109/sp54263.2024.00077}} \\
68 & I1 & Zufferey & 2024 & Our Data, Our Solutions: A Participatory Approach for Enhancing Privacy in Wearable Activity Tracker Third-Party Apps & \href{https://doi.org/10.56553/popets-2024-0139}{\nolinkurl{10.56553/popets-2024-0139}} \\
69 & I1 & Özçağdavul & 2024 & General Data Protection Regulation Compliance and Privacy Protection in Wearable Health Devices: Challenges and Solutions & \href{https://doi.org/10.58252/artukluhealth.1566573}{\nolinkurl{10.58252/artukluhealth.1566573}} \\
70 & I1 & Forsberg & 2025 & Security Analysis of Top-Ranked mHealth Fitness Apps: An Empirical Study & \href{https://doi.org/10.1007/978-3-031-79007-2_19}{\nolinkurl{10.1007/978-3-031-79007-2\_19}} \\
71 & I1 & Hudig & 2025 & Intimate Data Sharing: Enhancing Transparency and Control in Fertility Tracking & \href{https://doi.org/10.1145/3706598.3714089}{\nolinkurl{10.1145/3706598.3714089}} \\
72 & I1 & Ikegwu & 2025 & Cyber threats in mobile healthcare applications: systematic review of enabling technologies, threat models, detection approaches, and future directions & \href{https://doi.org/10.1007/s10791-025-09686-z}{\nolinkurl{10.1007/s10791-025-09686-z}} \\
73 & I1 & Lange & 2025 & Slice it up: Unmasking User Identities in Smartwatch Health Data & \href{https://doi.org/10.1145/3708821.3710834}{\nolinkurl{10.1145/3708821.3710834}} \\
74 & I1 & Kenderova & 2026 & “It’s Always in the Back of My Mind”: Navigating sensitive health data amid reduced abortion rights in the United States & \href{https://doi.org/10.1145/3788682}{\nolinkurl{10.1145/3788682}} \\
75 & I1 & Nickel & 2026 & The illusion of empowerment: commercializing women’s health in the digital age & \href{https://doi.org/10.1093/heapro/daag074}{\nolinkurl{10.1093/heapro/daag074}} \\
76 & I1 & Pybus & 2026 & Extraction-by-design: Auditing infrastructures of datafication in baby-tracking apps & \href{https://doi.org/10.14763/2026.1.2087}{\nolinkurl{10.14763/2026.1.2087}} \\
77 & I2 & Paul & 2014 & Privacy Implications of Wearable Health Devices & \href{https://doi.org/10.1145/2659651.2659683}{\nolinkurl{10.1145/2659651.2659683}} \\
78 & I2 & Sunyaev & 2014 & Availability and quality of mobile health app privacy policies & \href{https://doi.org/10.1136/amiajnl-2013-002605}{\nolinkurl{10.1136/amiajnl-2013-002605}} \\
79 & I2 & Huckvale & 2015 & Unaddressed privacy risks in accredited health and wellness apps: a cross-sectional systematic assessment & \href{https://doi.org/10.1186/s12916-015-0444-y}{\nolinkurl{10.1186/s12916-015-0444-y}} \\
80 & I2 & Blenner & 2016 & Privacy Policies of Android Diabetes Apps and Sharing of Health Information. & \href{https://doi.org/10.1001/jama.2015.19426}{\nolinkurl{10.1001/jama.2015.19426}} \\
81 & I2 & Hilts & 2016 & Every step you fake: a comparative analysis of fitness tracker privacy and security & \href{https://apo.org.au/sites/default/files/resource-files/2016-02/apo-nid61229.pdf}{apo.org.au} \\
82 & I2 & Rosenfeld & 2017 & Data Security and Privacy in Apps for Dementia: An Analysis of Existing Privacy Policies & \href{https://doi.org/10.1016/j.jagp.2017.04.009}{\nolinkurl{10.1016/j.jagp.2017.04.009}} \\
83 & I2 & Albrecht & 2018 & Relevance of Trust Marks and CE Labels in German-Language Store Descriptions of Health Apps: Analysis & \href{https://doi.org/10.2196/10394}{\nolinkurl{10.2196/10394}} \\
84 & I2 & Hutton & 2018 & Assessing the Privacy of mHealth Apps for Self-Tracking: Heuristic Evaluation Approach & \href{https://doi.org/10.2196/mhealth.9217}{\nolinkurl{10.2196/mhealth.9217}} \\
85 & I2 & O'Loughlin & 2018 & Reviewing the data security and privacy policies of mobile apps for depression & \href{https://doi.org/10.1016/j.invent.2018.12.001}{\nolinkurl{10.1016/j.invent.2018.12.001}} \\
86 & I2 & Powell & 2018 & The Complexity of Mental Health App Privacy Policies: A Potential Barrier to Privacy & \href{https://doi.org/10.2196/mhealth.9871}{\nolinkurl{10.2196/mhealth.9871}} \\
87 & I2 & Schairer & 2018 & How Could Commercial Terms of Use and Privacy Policies Undermine Informed Consent in the Age of Mobile Health? & \href{https://doi.org/10.1001/amajethics.2018.864}{\nolinkurl{10.1001/amajethics.2018.864}} \\
88 & I2 & Huckvale & 2019 & Assessment of the Data Sharing and Privacy Practices of Smartphone Apps for Depression and Smoking Cessation & \href{https://doi.org/10.1001/jamanetworkopen.2019.2542}{\nolinkurl{10.1001/jamanetworkopen.2019.2542}} \\
89 & I2 & Mulder & 2019 & Privacy policies, cross-border health data and the GDPR & \href{https://doi.org/10.1080/13600834.2019.1644068}{\nolinkurl{10.1080/13600834.2019.1644068}} \\
90 & I2 & Mulder & 2019 & Health Apps, their Privacy Policies and the GDPR & \href{https://ejlt.org/index.php/ejlt/article/view/667}{ejlt.org} \\
91 & I2 & Parker & 2019 & How private is your mental health app data? An empirical study of mental health app privacy policies and practices & \href{https://doi.org/10.1016/j.ijlp.2019.04.002}{\nolinkurl{10.1016/j.ijlp.2019.04.002}} \\
92 & I2 & Robillard & 2019 & Availability, readability, and content of privacy policies and terms of agreements of mental health apps & \href{https://doi.org/10.1016/j.invent.2019.100243}{\nolinkurl{10.1016/j.invent.2019.100243}} \\
93 & I2 & Benjumea & 2020 & Privacy Assessment in Mobile Health Apps: Scoping Review & \href{https://doi.org/10.2196/18868}{\nolinkurl{10.2196/18868}} \\
94 & I2 & Benjumea & 2020 & Assessment of the Fairness of Privacy Policies of Mobile Health Apps: Scale Development and Evaluation in Cancer Apps & \href{https://doi.org/10.2196/17134}{\nolinkurl{10.2196/17134}} \\
95 & I2 & Fan & 2020 & An Empirical Evaluation of GDPR Compliance Violations in Android mHealth Apps & \href{https://doi.org/10.1109/issre5003.2020.00032}{\nolinkurl{10.1109/issre5003.2020.00032}} \\
96 & I2 & Melcher & 2020 & Smartphone Apps for College Mental Health: A Concern for Privacy and Quality of Current Offerings & \href{https://doi.org/10.1176/appi.ps.202000098}{\nolinkurl{10.1176/appi.ps.202000098}} \\
97 & I2 & Ni & 2020 & Privacy Policy Compliance of Chronic Disease Management Apps in China: Scale Development and Content Evaluation & \href{https://doi.org/10.2196/23409}{\nolinkurl{10.2196/23409}} \\
98 & I2 & Shipp & 2020 & How private is your period?: A systematic analysis of menstrual app privacy policies & \href{https://doi.org/10.2478/popets-2020-0083}{\nolinkurl{10.2478/popets-2020-0083}} \\
99 & I2 & Jilka & 2021 & Terms and conditions apply: Critical issues for readability and jargon in mental health depression apps & \href{https://doi.org/10.1016/j.invent.2021.100433}{\nolinkurl{10.1016/j.invent.2021.100433}} \\
100 & I2 & LaMonica & 2021 & Privacy Practices of Health Information Technologies: Privacy Policy Risk Assessment Study and Proposed Guidelines & \href{https://doi.org/10.2196/26317}{\nolinkurl{10.2196/26317}} \\
101 & I2 & Mehrnezhad & 2021 & Caring for Intimate Data in Fertility Technologies & \href{https://doi.org/10.1145/3411764.3445132}{\nolinkurl{10.1145/3411764.3445132}} \\
102 & I2 & Tangari & 2021 & Mobile health and privacy: cross sectional study & \href{https://doi.org/10.1136/bmj.n1248}{\nolinkurl{10.1136/bmj.n1248}} \\
103 & I2 & Alfawzan & 2022 & Privacy, Data Sharing, and Data Security Policies of Women’s mHealth Apps: Scoping Review and Content Analysis & \href{https://doi.org/10.2196/33735}{\nolinkurl{10.2196/33735}} \\
104 & I2 & Camacho & 2022 & Assessment of Mental Health Services Available Through Smartphone Apps & \href{https://doi.org/10.1001/jamanetworkopen.2022.48784}{\nolinkurl{10.1001/jamanetworkopen.2022.48784}} \\
105 & I2 & Dong & 2022 & Privacy Analysis of Period Tracking Mobile Apps in the Post-Roe v. Wade Era & \href{https://doi.org/10.1145/3551349.3561343}{\nolinkurl{10.1145/3551349.3561343}} \\
106 & I2 & Erickson & 2022 & What You Do Not Expect When You Are Expecting: Privacy Analysis of Femtech & \href{https://doi.org/10.1109/tts.2022.3160928}{\nolinkurl{10.1109/tts.2022.3160928}} \\
107 & I2 & Echenim & 2023 & Ensuring Privacy Policy Compliance of Wearables with IoT Regulations & \href{https://doi.org/10.1109/tps-isa58951.2023.00039}{\nolinkurl{10.1109/tps-isa58951.2023.00039}} \\
108 & I2 & Neal & 2023 & Read and accepted? Scoping the cognitive accessibility of privacy policies of health apps and websites in three European countries & \href{https://doi.org/10.1177/20552076231152162}{\nolinkurl{10.1177/20552076231152162}} \\
109 & I2 & Singh & 2023 & Assessment of App Store Description and Privacy Policy to Explore Ethical and Safety Concerns Associated with the Use of Mental Health Apps for Depression & \href{https://doi.org/10.1177/02537176221142046}{\nolinkurl{10.1177/02537176221142046}} \\
110 & I2 & Krumbholz & 2024 & Tracking the Trackers: ‘Menstruapp’ Privacy Policies Following the Dobbs Decision & \href{https://doi.org/10.1080/1554477x.2024.2300103}{\nolinkurl{10.1080/1554477x.2024.2300103}} \\
111 & I2 & Malki & 2024 & Exploring Privacy Practices of Female mHealth Apps in a Post-Roe World & \href{https://doi.org/10.1145/3613904.3642521}{\nolinkurl{10.1145/3613904.3642521}} \\
112 & I2 & Saini & 2024 & Privacy and Security of Women's Reproductive Health Apps in a Changing Legal Landscape & \href{https://doi.org/10.48550/arxiv.2404.05876}{\nolinkurl{10.48550/arxiv.2404.05876}} \\
113 & I2 & Zadushlivy & 2024 & Exploration of Reproductive Health Apps’ Data Privacy Policies and the Risks Posed to Users: Qualitative Content Analysis & \href{https://doi.org/10.2196/51517}{\nolinkurl{10.2196/51517}} \\
114 & I2 & Doherty & 2025 & Privacy in consumer wearable technologies: a living systematic analysis of data policies across leading manufacturers & \href{https://doi.org/10.1038/s41746-025-01757-1}{\nolinkurl{10.1038/s41746-025-01757-1}} \\
115 & I2 & Grimme & 2025 & Hidden in Plain Sight: a Structured Analysis of Privacy Policies in the Context of Body-worn 'FemTech' Technologies & \href{https://doi.org/10.1145/3706598.3713702}{\nolinkurl{10.1145/3706598.3713702}} \\
116 & I2 & Hassan & 2025 & Unveiling Privacy and Security Gaps in Female Health Apps & \href{https://doi.org/10.48550/arxiv.2502.02749}{\nolinkurl{10.48550/arxiv.2502.02749}} \\
117 & I2 & Javed & 2025 & Evaluating privacy policies of AI-powered mHealth iOS applications & \href{https://doi.org/10.1093/jamia/ocaf130}{\nolinkurl{10.1093/jamia/ocaf130}} \\
118 & I2 & Lin & 2025 & Quality and Privacy Policy Compliance of Mental Health Care Apps in China: Cross-Sectional Evaluation Study & \href{https://doi.org/10.2196/66762}{\nolinkurl{10.2196/66762}} \\
119 & I2 & Stevenson & 2025 & “Your Doctor is Spying on You”: An Analysis of Data Practices in Mobile Healthcare Applications & \href{https://doi.org/10.1109/buildsec68439.2025.00021}{\nolinkurl{10.1109/buildsec68439.2025.00021}} \\
120 & I2 & Tzanou & 2025 & Research Report: Femtech Apps and Wearables and GDPR Compliance- An Empirical Study of Femtech's Privacy Policies and Terms of Service under the GDPR & \href{https://doi.org/10.2139/ssrn.5177016}{\nolinkurl{10.2139/ssrn.5177016}} \\
121 & I2 & White & 2025 & Experiential case study audit of three popular period trackers using General Data Protection Regulation (GDPR) and intimate privacy assessment criteria & \href{https://doi.org/10.1016/j.contraception.2025.111235}{\nolinkurl{10.1016/j.contraception.2025.111235}} \\
122 & I2 & Graham & 2026 & The need for transparency in the promotion of popular period tracker applications (apps): a content analysis of their app store descriptions & \href{https://doi.org/10.1080/14647273.2025.2607988}{\nolinkurl{10.1080/14647273.2025.2607988}} \\
123 & I3 & Aamir & 2023 & GenAIPABench: A Benchmark for Generative AI-based Privacy Assistants & \href{https://doi.org/10.13016/m29yec-3iy3}{\nolinkurl{10.13016/m29yec-3iy3}} \\
124 & I3 & Tang & 2023 & PolicyGPT: Automated Analysis of Privacy Policies with Large Language Models & \href{https://doi.org/10.48550/arxiv.2309.10238}{\nolinkurl{10.48550/arxiv.2309.10238}} \\
125 & I3 & Bolun & 2024 & Empowering Users in Digital Privacy Management through Interactive LLM-Based Agents & \href{https://doi.org/10.48550/arxiv.2410.11906}{\nolinkurl{10.48550/arxiv.2410.11906}} \\
126 & I3 & Göknil & 2024 & Privacy Policy Analysis through Prompt Engineering for LLMs & \href{https://doi.org/10.48550/arxiv.2409.14879}{\nolinkurl{10.48550/arxiv.2409.14879}} \\
127 & I3 & Morales & 2024 & A Large Language Model Approach to Code and Privacy Policy Alignment & \href{https://doi.org/10.1109/saner60148.2024.00016}{\nolinkurl{10.1109/saner60148.2024.00016}} \\
128 & I3 & Nguyen & 2024 & XLMR4MD: New Vietnamese dataset and framework for detecting the consistency of description and permission in Android applications using large language models & \href{https://doi.org/10.1016/j.cose.2024.103814}{\nolinkurl{10.1016/j.cose.2024.103814}} \\
129 & I3 & Rodríguez & 2024 & Large language models: a new approach for privacy policy analysis at scale & \href{https://doi.org/10.1007/s00607-024-01331-9}{\nolinkurl{10.1007/s00607-024-01331-9}} \\
130 & I3 & Silva & 2024 & Entailment-Driven Privacy Policy Classification with LLMs & \href{https://doi.org/10.1109/buildsec64048.2024.00010}{\nolinkurl{10.1109/buildsec64048.2024.00010}} \\
131 & I3 & Alecci & 2025 & Toward LLM-Driven GDPR Compliance Checking for Android Apps & \href{https://doi.org/10.1145/3696630.3728508}{\nolinkurl{10.1145/3696630.3728508}} \\
132 & I3 & Chanenson & 2025 & Automating Governing Knowledge Commons and Contextual Integrity (GKC-CI) Privacy Policy Annotations with Large Language Models & \href{https://doi.org/10.56553/popets-2025-0062}{\nolinkurl{10.56553/popets-2025-0062}} \\
133 & I3 & Xie & 2025 & Evaluating Privacy Policies under Modern Privacy Laws At Scale: An LLM-Based Automated Approach & \href{https://www.semanticscholar.org/paper/18b48b22e55758f15c1373057bb10a7981450e92}{semanticscholar.org} \\
134 & I3 & Yuxin & 2025 & Using LLMs for Automated Privacy Policy Analysis: Prompt Engineering, Fine-Tuning and Explainability & \href{https://doi.org/10.48550/arxiv.2503.16516}{\nolinkurl{10.48550/arxiv.2503.16516}} \\
135 & I3 & Cory & 2026 & Word-level Annotation of GDPR Transparency Compliance in Privacy Policies using Large Language Models & \href{https://doi.org/10.56553/popets-2026-0026}{\nolinkurl{10.56553/popets-2026-0026}} \\
136 & I4 & Reidenberg & 2014 & Privacy Harms and the Effectiveness of the Notice and Choice Framework & \href{https://doi.org/10.2139/ssrn.2418247}{\nolinkurl{10.2139/ssrn.2418247}} \\
137 & I4 & Bhatia & 2015 & Towards an information type lexicon for privacy policies & \href{https://doi.org/10.1109/relaw.2015.7330207}{\nolinkurl{10.1109/relaw.2015.7330207}} \\
138 & I4 & Schaub & 2015 & A Design Space for Effective Privacy Notices & \href{https://doi.org/10.1017/9781316831960.021}{\nolinkurl{10.1017/9781316831960.021}} \\
139 & I4 & Bhatia & 2016 & A Theory of Vagueness and Privacy Risk Perception & \href{https://doi.org/10.1109/re.2016.20}{\nolinkurl{10.1109/re.2016.20}} \\
140 & I4 & Gluck & 2016 & How Short Is Too Short? Implications of Length and Framing on the Effectiveness of Privacy Notices & \href{https://www.usenix.org/system/files/conference/soups2016/soups2016-paper-gluck.pdf}{usenix.org} \\
141 & I4 & Reidenberg & 2016 & Ambiguity in Privacy Policies and the Impact of Regulation & \href{https://doi.org/10.1086/688669}{\nolinkurl{10.1086/688669}} \\
142 & I4 & Strahilevitz & 2016 & Is Privacy Policy Language Irrelevant to Consumers? & \href{https://doi.org/10.1086/689933}{\nolinkurl{10.1086/689933}} \\
143 & I4 & Waldman & 2016 & Privacy, Notice, and Design & \href{https://doi.org/10.2139/ssrn.2780305}{\nolinkurl{10.2139/ssrn.2780305}} \\
144 & I4 & Chua & 2017 & Compliance to personal data protection principles: A study of how organizations frame privacy policy notices & \href{https://doi.org/10.1016/j.tele.2017.01.008}{\nolinkurl{10.1016/j.tele.2017.01.008}} \\
145 & I4 & Lebanoff & 2018 & Automatic Detection of Vague Words and Sentences in Privacy Policies & \href{https://doi.org/10.18653/v1/d18-1387}{\nolinkurl{10.18653/v1/d18-1387}} \\
146 & I4 & Liu & 2018 & Modeling Language Vagueness in Privacy Policies using Deep Neural Networks & \href{https://doi.org/10.48550/arxiv.1805.10393}{\nolinkurl{10.48550/arxiv.1805.10393}} \\
147 & I4 & Shvartzshnaider & 2019 & Going against the (Appropriate) Flow: A Contextual Integrity Approach to Privacy Policy Analysis & \href{https://doi.org/10.1609/hcomp.v7i1.5266}{\nolinkurl{10.1609/hcomp.v7i1.5266}} \\
148 & I4 & Zou & 2019 & You `Might' Be Affected & \href{https://doi.org/10.1145/3290605.3300424}{\nolinkurl{10.1145/3290605.3300424}} \\
149 & I4 & Ebert & 2020 & Does Context in Privacy Communication Really Matter? A Survey on Consumer Concerns and Preferences & \href{https://doi.org/10.1145/3313831.3376575}{\nolinkurl{10.1145/3313831.3376575}} \\
150 & I4 & Samat & 2020 & Format vs. Content: The impact of risk and presentation on disclosure decisions & \href{https://doi.org/10.1184/r1/13090571}{\nolinkurl{10.1184/r1/13090571}} \\
151 & I4 & Zhang & 2020 & COVID-19 Contact-Tracing Apps: Analysis of the Readability of Privacy Policies & \href{https://doi.org/10.2196/21572}{\nolinkurl{10.2196/21572}} \\
152 & I4 & Chen & 2021 & Fighting the Fog & \href{https://doi.org/10.1145/3463676.3485601}{\nolinkurl{10.1145/3463676.3485601}} \\
153 & I4 & Ebert & 2021 & Bolder is Better: Raising User Awareness through Salient and Concise Privacy Notices & \href{https://doi.org/10.1145/3411764.3445516}{\nolinkurl{10.1145/3411764.3445516}} \\
154 & I4 & Egliston & 2021 & Examining visions of surveillance in Oculus’ data and privacy policies, 2014–2020 & \href{https://doi.org/10.1177/1329878x211041670}{\nolinkurl{10.1177/1329878x211041670}} \\
155 & I4 & Li & 2022 & Understanding iOS Privacy Nutrition Labels: An Exploratory Large-Scale Analysis of App Store Data & \href{https://doi.org/10.1145/3491101.3519739}{\nolinkurl{10.1145/3491101.3519739}} \\
156 & I4 & Belcheva & 2023 & Understanding Website Privacy Policies—A Longitudinal Analysis Using Natural Language Processing & \href{https://doi.org/10.3390/info14110622}{\nolinkurl{10.3390/info14110622}} \\
157 & I4 & Brunotte & 2023 & Context, Content, Consent - How to Design User-Centered Privacy Explanations (S) & \href{https://doi.org/10.18293/seke2023-032}{\nolinkurl{10.18293/seke2023-032}} \\
158 & I4 & Kyi & 2023 & “It doesn’t tell me anything about how my data is used”: User Perceptions of Data Collection Purposes & \href{https://doi.org/10.1145/3613904.3642260}{\nolinkurl{10.1145/3613904.3642260}} \\
159 & I4 & Lian & 2023 & Really Vague? Automatically Identify the Potential False Vagueness within the Context of Documents & \href{https://doi.org/10.3390/math11102334}{\nolinkurl{10.3390/math11102334}} \\
160 & I4 & Cirucci & 2024 & Oversharing the super safe stuff: “Privacy-washing” in Apple iPhone and Google Pixel commercials & \href{https://doi.org/10.5210/fm.v29i5.13321}{\nolinkurl{10.5210/fm.v29i5.13321}} \\
161 & I4 & Ghahremani & 2024 & Comprehensive evaluation of privacy policies using the contextual integrity framework & \href{https://doi.org/10.1002/spy2.380}{\nolinkurl{10.1002/spy2.380}} \\
162 & I4 & Green & 2024 & Strategic Indeterminacy and Online Privacy Policies: (Un)informed Consent and the General Data Protection Regulation & \href{https://doi.org/10.1007/s11196-024-10132-4}{\nolinkurl{10.1007/s11196-024-10132-4}} \\
163 & I4 & Korneeva & 2024 & How Users Assess Privacy Risks in the Internet of Things: The Role of Framing, Comparing, and Educating & \href{https://doi.org/10.1177/00076503241255082}{\nolinkurl{10.1177/00076503241255082}} \\
164 & I4 & Lin & 2024 & Data Safety vs. App Privacy: Comparing the Usability of Android and iOS Privacy Labels & \href{https://doi.org/10.56553/popets-2024-0047}{\nolinkurl{10.56553/popets-2024-0047}} \\
165 & I4 & Pan & 2024 & Read or skip privacy policies when installing apps on wearable devices: the roles of perceived necessity and threat clues & \href{https://doi.org/10.1057/s41599-024-02989-4}{\nolinkurl{10.1057/s41599-024-02989-4}} \\
166 & I4 & Seiling & 2024 & Beware: Processing of Personal Data—Informed Consent Through Risk Communication & \href{https://doi.org/10.1109/tpc.2024.3361328}{\nolinkurl{10.1109/tpc.2024.3361328}} \\
167 & I4 & Chen & 2025 & Governing data through privacy after China's PIPL: Discursive strategies, layered compliance and power dynamics in digital platform policies & \href{https://doi.org/10.1016/j.telpol.2025.103147}{\nolinkurl{10.1016/j.telpol.2025.103147}} \\
168 & I4 & Romero & 2025 & The power of words: how linguistic framing affects consent in retail privacy policies & \href{https://doi.org/10.1016/j.jretai.2025.09.001}{\nolinkurl{10.1016/j.jretai.2025.09.001}} \\


\end{longtable}
}



\cleardoublepage
\pagestyle{fancy-lists}

% AI Appendix of the RUB template (bossthesis): unnumbered, after the lettered
% appendices. The declaration box at the front refers to it.
\chapter*{AI Appendix}
\addcontentsline{toc}{chapter}{AI Appendix}
\chaptermark{AI Appendix}

This appendix records the use of AI tools for the conceptual design of the
thesis and for its implementation, beyond the list in the declaration.

\section*{AI Use for the Concept}
\todo[inline]{State which AI tools were used for brainstorming or for the design and
conceptual aspects of the thesis, and to what extent; separate own ideas from
AI input.}

\section*{AI Use for the Implementation}
\todo[inline]{State which AI tools were used for which programs and scripts, and to
which end (for example debugging or code generation). The repository's commit
history records Claude Code as co-author of the scripts in \texttt{scripts/}.}



\end{document}
